\begin{abstract}
We give a deterministic integer-factorization algorithm with bit complexity
\(\exp(O_\varepsilon((\log N)^{2/3+\varepsilon}))\) for every fixed
\(\varepsilon>0\).  The construction modifies the exponent-annihilation
method in the pinned Cogentic manuscript on deterministic factoring.  Its
main ingredients are preconditioned polynomial pairs, squarefree smooth
values in arithmetic progressions, exact cancellation in two real-unit
lattices, and a supplied-exponent splitting algorithm.  We give an explicit
replacement parameter schedule and prove the uniform analytic transfer,
including the conditional Gaussian laws, exceptional-conductor estimates,
adaptive deletion argument, and final spectral lower bound.  A corrected
Perron estimate retains a term omitted from the source's displayed bound.
The algorithm uses finite enumeration and exact certificates; probability
is used to prove that a successful enumerated pair exists.  We also identify
the precise large-modulus obstruction in the present progression argument.
The result is a mathematical manuscript using stated classical inputs; it
does not assert verification of the source's formal-proof artifacts or an
endpoint \(L_N[2/3,c]\) bound.
\end{abstract}

\tableofcontents

\section{The result, its provenance, and the proof strategy}
\label{sec:introduction}

All logarithms without a subscript are natural.  Complexity means bit
operations on a deterministic multitape Turing machine.  An integer
\(N\ge4\) is supplied in binary.  Write \(L=\log N\), \(L_2=\log L\), and
\(\kappa=\log2\).  A returned factor is always checked to be an integer
\(g\) with \(1<g<N\) and \(g\mid N\).

\begin{theorem}\label{thm:main}
For every fixed rational \(0<\delta\le1/100\), there is an explicit
deterministic algorithm that factors every composite integer \(N\ge4\)
within
\[
 \exp\bigl(O_\delta(L^{h})\bigr),
 \qquad h=\frac23+8\delta,
\]
bit operations.  The constants are effective.  Consequently, for every
fixed \(\varepsilon>0\), deterministic factoring is possible in
\(\exp(O_\varepsilon(L^{2/3+\varepsilon}))\) bit operations.
\end{theorem}

Here ``factors'' can mean either finding one proper nontrivial factor or
producing the complete prime factorization.  The latter follows by
recursion and deterministic primality testing; at most \(O(\log N)\)
nontrivial recursive calls are needed, which is absorbed by the bound.
One may use AKS primality testing \cite{AKS2004}.  Perfect powers are
recognized by exact integer-root computations for exponents at most
\(\lfloor\log_2N\rfloor\).

The conventional notation
\[
 L_N[\alpha,c]=
 \exp\bigl((c+o(1))L^\alpha L_2^{1-\alpha}\bigr)
\]
should be kept separate from a fixed-power estimate.  The theorem supplies
a family of bounds approaching power \(2/3\); it does not supply a fixed
constant \(c\) in an endpoint \(L_N[2/3,c]\) estimate.  In particular, it
does not establish an algorithm with first \(L\)-parameter below \(1/3\).

\subsection{What is inherited and what is proved here}

The starting point is the deterministic Cogentic manuscript
\cite{CogDet2026}, frozen at Google Research commit
\texttt{4a7bc8ab0bc770e71a259b7966c500496b99010a}, file
\texttt{cogentic\_research\_papers/deterministic\_factoring\_l34/paper.pdf}.
Its Git blob is
\texttt{973002e0eb92acf3964fc020ca52ba6f1ac384ac}.
References to source lemma numbers below always mean that version.
The source proves a bound with power \(3/4\).  We retain its finite
algebraic architecture and replace its decimal parameter powers.

The new work is the uniform transfer, not merely an optimization of a
list of inequalities.  We check the progression estimate at the actual
moduli; the lattice congruences underlying the Gaussian conditioning;
the centered moment identity; the exceptional-conductor contributions;
the deletion process with its dependence on previously deleted rows;
and the spectral projection needed by the exact extractor.  The
parameter-free algebraic statements that we import are given with their
hypotheses in Section~\ref{sec:extraction}.  The supplied-exponent splitter
is stated in Section~\ref{sec:splitter}.  Its proof is developed in the
companion order-splitting manuscript, including a direct replacement for
the source's biquadratic enumeration implementation.

\begin{remark}[Status and reproducibility]
This is an ordinary mathematical proof using the explicitly stated
analytic and algebraic inputs.  The proof does not infer an asymptotic
running time from numerical experiments.  The accompanying exact
parameter checks are aids to reconstruction.  The source's full Lean
formalization has not been checked as part of this manuscript.  Pinning
the source fixes the object being used; it is not itself a mathematical
correctness certificate.
\end{remark}

\subsection{The algorithm in one paragraph}

For many primes \(\ell\in(B/2,B]\), construct two degree-\(d\) polynomials
that vanish at a common integer \(m\) modulo \(N\).  The first has \(d\)
simple roots modulo \(\ell\), and the second has none.  Enumerate small
linear perturbations preserving these facts.  For a successful pair,
jointly smooth polynomial values provide enough relations to cancel
all non-\(\ell\) ideal valuations and all real-unit logarithms.  Exact
rational reconstruction then gives an even \(E_\ell\) with
\(\ell^{E_\ell}=1\pmod N\).  A combinatorial saturation argument converts
enough such certificates into a common multiple of every \(p-1\),
\(p\mid N\).  A deterministic order splitter then returns a factor.

The analytic task is to show that the finite perturbation box contains
a successful pair for sufficiently many \(\ell\).  We prove this using
independent discrete Gaussians on that box's ambient lattice.  The actual
algorithm exhausts the box in a fixed order.

\begin{center}
\begin{tabular}{p{.25\textwidth}p{.33\textwidth}p{.32\textwidth}}
\toprule
Component & Required conclusion & Location \\
\midrule
Base-polynomial construction & Primitive, irreducible, real-rooted pairs;
fixed modular behavior & Sections~\ref{sec:parameters}--\ref{sec:construction} \\
Smooth progressions & Uniform squarefree smooth density outside controlled
conductors & Sections~\ref{sec:saddle}--\ref{sec:progressions} \\
Conditioned moments & Exact local factors and a centered square-moment bound
& Sections~\ref{sec:gaussian}--\ref{sec:moments} \\
Peeling and projection & A surviving positive Gram matrix with gap
at least \(L^{-3}\) & Sections~\ref{sec:peeling}--\ref{sec:spectral} \\
Exact extraction & Verified order multiples from that spanning certificate
& Section~\ref{sec:extraction} \\
Saturation and splitting & A common \(p-1\) multiple and a factor
& Sections~\ref{sec:saturation}--\ref{sec:runtime} \\
\bottomrule
\end{tabular}
\end{center}

\section{An explicit replacement parameter schedule}
\label{sec:parameters}

Fix \(0<\delta\le1/100\), rational, and set
\begin{align*}
 a&=\frac13,& t&=\frac13+\delta,& k&=\frac13+2\delta,
 &v&=\frac13+\frac52\delta,\\
 b&=\frac13+3\delta,& \eta&=\frac13+4\delta,
 &r&=\frac13+5\delta,& c&=\frac13+6\delta,\\
 h&=\frac23+8\delta,& m_0&=\frac23,
 &g&=1-b+\delta=\frac23-2\delta.
\end{align*}
The letter \(r\) denotes a power in this display; a modular root will
always be identified by its surrounding notation.  The roles of the
parameters are as follows.
\begin{center}
\begin{tabular}{lll}
\toprule
Quantity & Definition up to bounded integer rounding & Leading scale \\
\midrule
Degree, small-prime cutoff & \(d=z_0\) & \(L^a\) \\
Root-window width & \(w=2^{-L^t+O(1)}\) & \(\log(1/w)\asymp L^t\) \\
Approximation parameter & \(K_1=2^{L^k+O(1)}\) & \(\log K_1\asymp L^k\) \\
Smoothness bound & \(B=2^{L^b+O(1)}\) & \(\log B\asymp L^b\) \\
Evaluation range & \(A=2^{L^c+O(1)}\) & \(\log A\asymp L^c\) \\
Perturbation width & \(U=2^{L^h+O(1)}\) & \(\log U\asymp L^h\) \\
Base & \(m=\lceil N^{1/d}\rceil\) & \(\log m=(1+o(1))L^{m_0}\) \\
Polynomial value scale & \(X_{\rm sc}=c_dA^d\) & \(\log X_{\rm sc}=(\kappa+o(1))L^h\) \\
Smoothness parameter & \(u=\log X_{\rm sc}/\log B\) & \(u=(1+o(1))L^r\) \\
\bottomrule
\end{tabular}
\end{center}

\subsection{Computable rounding}

All powers in the algorithm are computed by rational arithmetic and
integer roots.  To be explicit, put \(n=\lfloor\log_2N\rfloor+1\),
\(z=1-N/2^n\in[0,1/2]\), and \(J=n+30\).  Set
\[
 \kappa_- =\sum_{j=1}^J\frac1{j2^j},\qquad
 \kappa_+=\kappa_-+2^{-J},\qquad
 \widehat L=n\kappa_+-\sum_{j=1}^J\frac{z^j}{j},\qquad
 L_0=\lceil\widehat L\rceil.
\]
The positive tails of the two convergent series show that
\(L<\widehat L<L+2^{-20}\) for \(N\ge16\), and
\(\kappa_-<\kappa<\kappa_+\).  For rational \(\theta=A_\theta/B_\theta\)
in \((0,1)\), let
\[
 e(\theta)=\max\{j\in\mathbb Z_{\ge0}:
       j^{B_\theta}\le \widehat L^{A_\theta}\}.
\]
The inequality is tested after clearing denominators.  Binary search
computes \(e(\theta)\) in polynomial time in \(\log N\), for fixed
\(\theta\), and \(e(\theta)=L^\theta+O(1)\).
Take \(d=z_0=e(a)\), \(B=2^{e(b)}\), \(K_1=2^{e(k)}\),
\(A=2^{e(c)}\), \(U=2^{e(h)}\), and \(w=2^{-e(t)}\).
Inputs for which a later finite guard fails go to trial division.

Once \(d\ge10\) (otherwise the algorithm uses fallback), put
\begin{gather*}
 P_0=\prod_{p\le z_0}p,\quad P_{\rm odd}=P_0/2,
 \quad S_0=UA,\quad S_1=U,\quad
 \bar u_j=L_0^{20}S_j\quad(j=0,1),\\
 \mathcal B_U=([1,2\bar u_0]\cap\mathbb Z)
                 \times([1,2\bar u_1]\cap\mathbb Z),
 \qquad \Delta_{\rm close}=AK_1^{-1/2},\\
 \mathcal P_{\rm gen}=\{\ell\text{ prime}:B/2<\ell\le B\}.
\end{gather*}
Choose the least positive integer \(j\) such that
\(3^jP_{\rm odd}A^d\ge6K_1BP_0m\bar u_0\), and define
\begin{equation}\label{eq:scale}
 c_d=P_{\rm odd}3^j,
 \qquad
 6K_1BP_0m\bar u_0\le X_{\rm sc}=c_dA^d
            <18K_1BP_0m\bar u_0.
\end{equation}
Finally define the integer saturation cutoff
\begin{equation}\label{eq:kidx}
 K_{\rm idx}=B\,2^{\left\lceil
       n\lceil\log_2(6L_0)\rceil/e(b)\right\rceil}.
\end{equation}
It satisfies
\begin{equation}\label{eq:kidx-lower}
 K_{\rm idx}\ge B\exp\left(\frac{L}{\log B}\log(6L)\right),
 \qquad
 \log K_{\rm idx}=O(L^{1-b}L_2).
\end{equation}

\subsection{Strict inequalities used in the proof}

The following collection records all the power comparisons that will
be used.  They are supplemented by the actual estimates in the ensuing
sections.
\begin{align}
 &a<t<k<v<b<\eta<r<c<m_0<h<3/4,
       \qquad a+c<h,\label{eq:gates1}\\
 &m_0<2b,\quad m_0-v<b,\quad m_0-b<k,
       \quad 2c<h+b,\label{eq:gates2}\\
 &\eta+\delta<c,\quad 1-b<g<h,
       \quad 2h+r=\frac53+21\delta<3.\label{eq:gates3}
\end{align}
The equality \(r=h-b\) is built into the schedule.  Integer rounding of
\(d\) gives the slightly more precise statement
\(\log m=L^{2/3}+O(L^{1/3})\), which fits all conductor caps below.

The two error scales for the sieve and the centered moments are
\begin{equation}\label{eq:errors}
 E_{\rm Br}=\exp\left(-\frac\delta2L^\delta L_2\right),
 \qquad E=E_{\rm rel}
 =\exp\left(-\frac{3\delta}4L^\delta L_2\right).
\end{equation}
Both are smaller than every fixed inverse power of \(L\).  However,
they are much larger than \(B^{-c_1}\) or \(K_1^{-c_1}\) for any fixed
\(c_1>0\).  This separation is used explicitly, especially in the
deletion argument.

\section{The finite polynomial construction and the deterministic search}
\label{sec:construction}

We work after removing even inputs and perfect powers.  Then
\((m-1)^d<N<m^d\).  Trial division up to \(K_{\rm idx}\) ensures that
\(P^-(N)>K_{\rm idx}\) if no factor has already been found.  Testing
\(\gcd(a-m,N)\) for \(1\le a\le A\) either gives a factor or guarantees
the localization condition used later.

For \(i\in\{1,2\}\) and \(1\le s\le d\), set
\[
 \rho_{i,s}=\frac{2s-2+i}{4d},\qquad
 P_i(t)=\prod_{s=1}^d(t-\rho_{i,s}).
\]
For each \(\ell\in\mathcal P_{\rm gen}\), let \(r_1<\cdots<r_d\)
be the first \(d\) residues modulo \(\ell\) different from \(m\bmod\ell\).
The base-polynomial construction of source Lemma~8.1 produces
\(f_i^{(0)}\in\mathbb Z[X]\) of degree \(d\), leading coefficient
\(c_d\), and integers \(1\le k_i<2c_d\), with
\begin{align}
 &f_i^{(0)}(m)=k_iN,\quad
       (k_1k_2,P_0\ell)=1,\label{eq:base-eval}\\
 &\left|[X^j]f_i^{(0)}-c_dA^{d-j}[t^j]P_i\right|
       <P_0\ell m\qquad(0\le j<d),\label{eq:base-approx}\\
 &f_1^{(0)}(X)\equiv c_d\prod_{s=1}^d(X-r_s)\pmod\ell,
       \qquad f_2^{(0)}\text{ has no root modulo }\ell.\label{eq:base-ell}
\end{align}
Every perturbation
\begin{equation}\label{eq:perturb}
 f_{i,\mathbf u}(X)
 =f_i^{(0)}(X)+P_0\ell(X-m)(u_0+u_1X),
 \qquad \mathbf u=(u_0,u_1)\in\mathbb Z^2,
\end{equation}
is primitive, remains irreducible modulo \(2\), and satisfies
\(f_{i,\mathbf u}(a)\equiv1\pmod{P_0}\) for every integer \(a\).
It is therefore irreducible over \(\mathbb Q\).

\paragraph{Why this construction is finite and effective.}
Choose the first irreducible monic degree-\(d\) polynomial over
\(\mathbb F_2\) by testing at most \(2^d\) candidates.  Modulo the odd
part of \(P_0\), prescribe the constant polynomial \(1\).  Modulo
\(\ell\), prescribe the split polynomial in \eqref{eq:base-ell} and a
root-free polynomial \(c_d(X^d-X-s_\ell)\).  Such an \(s_\ell\) exists
because \(x^d-x\) takes the same value at \(0\) and \(1\), hence misses
at least one residue.  Use the Chinese remainder theorem to round each
target coefficient to its prescribed class modulo \(P_0\ell\).
A balanced base-\(m\) correction of degree less than \(d\) enforces
\eqref{eq:base-eval} while preserving the congruences and
\eqref{eq:base-approx}; this is the coefficient-by-coefficient algorithm
in source Lemma~8.1.  Its time is polynomial in
\(2^d,\ell,z_0,\log c_d,d\log(mA)\), hence polynomial in \(B\).
Primitivity follows because every prime dividing \(c_d\) lies in
\(P_0\), while the constant term is \(1\) modulo each such prime.

Define
\begin{gather*}
 T_0=[3/4,3/4+w],\quad
 T_s=[\rho_{1,s}+w,\rho_{1,s}+2w],\quad
 T_{d+s}=[\rho_{2,s}+w,\rho_{2,s}+2w],\\
 I_j=AT_j\cap\mathbb Z,\qquad
 \mathcal A=\bigsqcup_{j=0}^{2d}I_j,\qquad
 \mathcal A^*=\{a\in\mathcal A:(a-m,k_1k_2)=1\},\\
 \mathcal J=\{p\text{ prime}:\sqrt B<p\le B,\ p\nmid k_1k_2\}.
\end{gather*}
All these sets are computed by trial division and gcd tests.  In
particular, \(\ell\in\mathcal J\), and no prime of \(\mathcal J\)
divides \(c_dN\).

\begin{lemma}[Uniform construction bounds]\label{lem:construction}
For all sufficiently large \(N\), uniformly in \(\ell\in\mathcal
P_{\rm gen}\) and \(\mathbf u,\mathbf v\in\mathcal B_U\), the two
perturbed polynomials satisfy:
\begin{enumerate}[label=(\roman*)]
\item \(\log c_d=(\kappa+o(1))L^h\),
\(\log X_{\rm sc}=(\kappa+o(1))L^h\), and
\(\|f_i\|_\infty<2X_{\rm sc}<2^{L_0}\).
\item On \([-1,1]\),
\(\left|X_{\rm sc}^{-1}f_i(At)-P_i(t)\right|\le K_1^{-1}\).
Each \(f_i\) has \(d\) distinct real roots \(\theta_{i,s}\), with
\(\left|\theta_{i,s}/A-\rho_{i,s}\right|<K_1^{-0.95}\).
\item For \(a\in\mathcal A\),
\[
 w(8d)^{-d}\le |P_i(a/A)|\le1,
 \quad
 \min_s|a/A-\theta_{i,s}/A|\ge w/2,
\]
and replacing the target roots by the actual roots changes each
logarithm by at most \(K_1^{-0.90}\).
\item A squarefree \(\mathcal J\)-smooth value has at most \(2.1u\)
prime factors.  Thus a joint relation has at most \(5u<L_0\) nonzero
valuation features, including the top row of ones.
\end{enumerate}
\end{lemma}

\begin{proof}
The logarithm of the right side of \eqref{eq:scale} is
\(\kappa L^h+L^{m_0}+O(L^c+L^k+L^b+L^a+L_2)\).
Since \(h>m_0\) and \(h>a+c\), subtracting \(d\log A\) gives the
first claim.  In particular, \(c_d\gg P_0B\), \(dA=o(m)\), and
\(dA^{d-1}=o(\bar u_0)\).  These are the explicit estimates used in
the coefficient-error and perturbation bounds of source Lemma~8.2(3).
Summing those errors on \([-1,1]\) and using \eqref{eq:scale} proves
the uniform \(K_1^{-1}\) approximation.

At distance \(K_1^{-0.95}\) on either side of a target root, the product
of the other \(d-1\) root distances is at least \((8d)^{-d}\).
Since \(\log K_1\gg d\log d\), the target polynomial's magnitude
there exceeds \(K_1^{-1}\); its two signs are opposite.  Degree \(d\)
then gives exactly one real root in each disjoint isolating interval.
The sampling intervals are distance at least \(w\) from their target
roots, and \(\log(1/w)=o(\log K_1)\), proving the logarithmic
stability.  Finally, each prime divisor of a smooth value is larger
than \(\sqrt B\), whereas \(\log|f_i(a)|\le(1+o(1))\log X_{\rm sc}\).
This gives the prime-factor count.  All comparisons are effective.
\end{proof}

\begin{algorithm}[Deterministic annihilator search]\label{alg:annihilator}
On an odd composite non-perfect-power input, perform the following
steps.  Every failure mentioned below invokes complete trial division.
\begin{enumerate}
\item For \(N\le64\), use fallback immediately.  Compute
\(\widehat L\) and \(d\) first; if \(d<10\), use fallback before
forming \(P_{\rm odd}\) or \(c_d\).  Compute the remaining rounded
parameters.  Trial-divide up to
\(\min(K_{\rm idx},N-1)\).  Test \(\gcd(a-m,N)\) for
\(1\le a\le A\); return a proper gcd and invoke fallback if a gcd is
\(N\).
\item Check the construction and extractor guards: \(d\ge10\), \(L_0\ge10\),
\(B\ge20\), \(B>z_0^2\), \(2(d+1)<B\), \(16dw\le1\),
\(4dA\le m\le2^{L_0}\), \(K_1\ge16\), \(4P_0B\le c_d\), and
\(\lceil\log_2K_1\rceil\le64d^3(d+1)L_0\).  Also check that the
root-isolating radius \(2^{-\lceil0.95\log_2K_1\rceil}\) is less
than \(1/(8d)\).
\item Initialize \(M_*=2\) and \(\mathcal P_{\rm ok}=\varnothing\).
For each \(\ell\in\mathcal P_{\rm gen}\), construct the base pair,
\(\mathcal J\), and \(\mathcal A^*\).  Enumerate
\((\mathbf u,\mathbf v)\in\mathcal B_U^2\) lexicographically.
When the polynomial heights are at most \(2^{L_0}\), call the exact
procedure \(\operatorname{TryPair}\) of
Lemma~\ref{lem:trypair}.  On its first successful certificate
\(E_\ell\), record \(\ell\), replace
\(M_*\) by \(\operatorname{lcm}(M_*,E_\ell)\), and proceed to the
next prime.
\item Require the rationally checkable inequality
\(6|\mathcal P_{\rm ok}|e(b)\kappa_-\ge B\).  Multiply \(M_*\)
by every maximal prime power at most \(K_{\rm idx}\), obtaining
\(M\).  Abort this multiplication if its bit length exceeds
\(2^{e(g)}\).  Return \(M\) if the cap holds; otherwise use fallback.
\end{enumerate}
\end{algorithm}

The analytic proof will show that fallback is never reached above an
effective threshold depending only on \(\delta\).  The algorithm itself
does not need to know that threshold.

\section{Preconditioning and the elementary weighted sieve}
\label{sec:sieve}

The factor \(k_1k_2\) in \eqref{eq:base-eval} cannot be ignored.  It is
the reason for the restriction defining \(\mathcal A^*\).  Every prime
dividing \(k_1k_2\) exceeds \(z_0\), and
\(k_1k_2<4c_d^2\).  Since \(h<3/4\), for sufficiently large \(N\)
the number of its distinct prime divisors is less than \(L^{0.76}\).
Comparison with the smallest such primes, followed by the prime number
theorem, gives
\begin{equation}\label{eq:coprime-density}
 \prod_{p\mid k_1k_2}(1-p^{-1})>\frac14.
\end{equation}
Indeed, all the relevant smallest primes lie below \(L^{0.80}\), and
\(\sum_{L^{1/3}<p\le L^{0.80}}p^{-1}=\log(2.4)+o(1)<1\).
The sum of their squares is negligible, so the Euler product exceeds
\(1/4\).  This use of an asymptotic estimate is effective.

Define
\[
 \psi_0(q)=q\prod_{p\mid(q,k_1k_2)}(1-p^{-1}).
\]
The following is the version of source Lemma~8.2(2) needed here.

\begin{lemma}[Weighted progression count]\label{lem:brun}
For every sampling interval \(I_j\),
\(|I_j\cap\mathcal A^*|\ge Aw/5\).  Suppose that \(q\) is squarefree
and supported on primes at most \(\sqrt B\), that
\((r_q-m,q)=1\), that \((Q,qk_1k_2)=1\), and that \(qQ\le B^4\).
If \(|F(a)|\le F_0\) and
\[
 |F(a+1)-F(a)|\le \frac{L^{40}F_0}{Aw}
 \quad(a,a+1\in I_j),
\]
then, uniformly in \(r_Q\),
\begin{align*}
 &\frac1{|I_j\cap\mathcal A^*|}
 \sum_{\substack{a\in I_j\cap\mathcal A^*\\
      a\equiv r_q\ (q),\ a\equiv r_Q\ (Q)}}F(a)\\
 &\hspace{1cm}=
 \frac1{\psi_0(q)Q}\,
 \frac1{|I_j\cap\mathcal A^*|}
 \sum_{a\in I_j\cap\mathcal A^*}F(a)
 +O\left(\frac{F_0E_{\rm Br}}{\psi_0(q)Q}\right).
\end{align*}
\end{lemma}

\begin{proof}
Take the sieve threshold \(z_{\rm Br}=\exp(L^\eta)\) and truncation
order \(r_0=\lceil L^\delta\rceil\).  Primes dividing \(k_1k_2\)
above \(z_{\rm Br}\) remove at most
\(O(\omega(k_1k_2)(|I|/(qQz_{\rm Br})+1))\) points from any
subinterval \(I\) in the progression.  For the remaining primes,
Bonferroni's inequalities at order \(2r_0\) bound the relative omitted
tail by
\[
 O\left(\sum_{\nu\ge2r_0}\frac{C^\nu}{\nu!}\right)
 =\exp\bigl(-(2\delta+o(1))L^\delta L_2\bigr),
\]
because their reciprocal sum is bounded absolutely.  There are at most
\(\exp(O(L^\delta L_2))\) retained divisors, each at most
\(z_{\rm Br}^{2r_0}\).  Thus every progression modulus used has
logarithm
\[
 O(L^b+L^{\eta+\delta})=o(L^c).
\]
Since \(\log(Aw)=(\kappa+o(1))L^c\), every endpoint error is
negligible, even on subintervals of length
\(AwE_{\rm Br}L_0^{-50}\).  Dividing the sieved count by the corresponding
unrestricted count gives the local factor \(\psi_0(q)^{-1}Q^{-1}\).
The density assertion follows from \eqref{eq:coprime-density}.
Partition into \(\lceil L^{40}/E_{\rm Br}\rceil\) consecutive
subintervals, whose lengths still exceed that lower cutoff for large
\(N\).  The displayed Lipschitz bound makes the variation on each
piece \(O(F_0E_{\rm Br})\).  The spare exponent between
\(2\delta\) and \(\delta/2\) absorbs its polynomial losses and yields
the asserted error.
\end{proof}

\section{Classical analytic inputs and a corrected prime-sum estimate}
\label{sec:analytic-inputs}

We state the analytic inputs explicitly.  The proof uses effective
versions of the prime number theorem, the Landau--Page theorem and its
explicit-formula consequence, and the near-one zero-density estimate.
These are the inputs recorded in source Proposition~7.1.

\begin{enumerate}[label=(A\arabic*)]
\item\label{input:pnt} For an effective \(c>0\),
\(\vartheta(x)=x+O(x\exp(-c\sqrt{\log x}))\).
Its consequences used below include partial summation for prime powers
and \(\sum_{p\le x}p^{-1}=\log\log x+C+o(1)\).
See \cite{PNT}.
\item\label{input:page} There is an effective absolute \(c_{\rm LP}>0\)
such that the primitive Dirichlet \(L\)-functions of conductor at most
\(D\) are zero-free in the usual region
\(\Re s\ge1-c_{\rm LP}/\log(D(|\Im s|+4))\), with at most one
exceptional primitive real character and one real exceptional zero.
The explicit formula in that range gives the exceptional character the
main term \(-x^\beta/\beta\), when its zero \(\beta\) is present;
the derivative of this main term is nonpositive.  We use the effective
uniform error stated in source Proposition~7.1(2), at
\(D=T=\exp(c_{\rm LP}\log B/(800L_2))\).
See \cite{LandauPage}.
\item\label{input:density} For \(4/5\le\sigma\le1\), \(Q\ge1\),
and \(T\ge2\),
\begin{equation}\label{eq:density}
 \sum_{d\le Q}\ \sum_{\substack{\chi\bmod d\\\chi\ {\rm primitive}}}
 N(\sigma,T,\chi)
 \ll (Q^2T)^{2(1-\sigma)/\sigma}\log^{14}(QT+2).
\end{equation}
Here \(N(\sigma,T,\chi)\) counts zeros with real part at least
\(\sigma\) and imaginary part of absolute value at most \(T\), with
multiplicity.  This is Montgomery's near-one estimate
\cite{Montgomery1971}; its exact displayed form is also used in the
proof of Lemma~5.2 of Ren--Tsang \cite{RenTsang2005}.
\end{enumerate}

\begin{remark}[Bibliographic precision]
The source attributes the last estimate to Montgomery's Theorem~12.1,
whereas Ren--Tsang attribute it to Theorem~12.2.  We use the explicit
inequality \eqref{eq:density}.  No conclusion depends on resolving this
numbering discrepancy.  A direct check of the original book's numbering
is not claimed here.
\end{remark}

For clarity, we retain the following slightly weaker, but fully adequate,
prime-sum bound.  It corrects the omission of a horizontal-contour factor
in source Lemma~9.2(1).

\begin{lemma}[Corrected zero-free prime-sum estimate]\label{lem:prime-sum}
Let \(\chi\) be primitive of conductor \(d_0\ge1\), let
\(0<\nu\le1/8\), \(H\ge8\), and
\(1-\nu/2\le\sigma\le1\).  Suppose \(L(s,\chi)\) is zero-free on
\[
 \Re s\ge1-\nu,\qquad |\Im s-t|\le2H.
\]
For the trivial character of conductor one assume also \(|t|\ge3H\).
Put \(\mathcal L=\log(d_0(|t|+H+4)/\nu)\).  Uniformly for
\(\sqrt B\le x_1<x_2\le B\),
\begin{align}\label{eq:prime-sum}
 \left|\sum_{x_1<p\le x_2}\chi(p)p^{-\sigma-it}\log p\right|
 \ll{}&\frac{\mathcal L}{\nu}
     x_1^{1-\nu/2-\sigma}\log(H\log B)\notag\\
 &+\frac{x_2^{1-\sigma}}{H}
       \left((\log B)^2+\frac{\mathcal L}{\nu}\right)
       +B^{-1/10}.
\end{align}
\end{lemma}

\begin{proof}
The Borel--Carath\'eodory argument on the zero-free rectangle bounds
\(L'/L\) on its inner part by \(O(\mathcal L/\nu)\).  Apply truncated
Perron to the difference of the two interval endpoints and shift to
\(\Re z=1-\nu/2-\sigma\).  The quotient
\((x_2^z-x_1^z)/z\) is analytic at zero, so no artificial pole is
introduced there.  The vertical integral gives the first term.
The horizontal segments give the second term with
\(\mathcal L/\nu\) retained; the truncated Perron endpoint error is
covered by its \((\log B)^2\) term.  The prime-power terms are
\(O(B^{-1/10})\) in the indicated range.  Half-integer endpoint
approximations remove endpoint coincidences without changing the bound.
These are the steps of source Lemma~9.2(1), with its displayed
horizontal-contour contribution preserved.
\end{proof}

Set, from now on,
\begin{equation}\label{eq:conductor-parameters}
 \begin{split}
 Q_{\max}&=\exp(3L^{m_0}),\qquad T_*=2K_1^{7/2},\\
 D_{\rm small}&=\exp\left(\frac{c_{\rm LP}\log B}{800L_2}\right),
 \qquad \nu_1=\frac{400L_2}{\log B},\qquad \nu_0=L^{-v}.
 \end{split}
\end{equation}
In applications take \(H=L^{12}\).  The high-conductor factor
\(\mathcal L/\nu_0\) is at most
\(L^{m_0+v+o(1)}<L^{1.1}\).  At the saddle below,
\(B^{1-\alpha}=O(u\log u)\le L^{0.4}\) eventually.  Thus the
horizontal term is negligible.  The vertical term contains
\[
 \exp(-c\nu_0\log B)
   =\exp(-cL^{b-v})=\exp(-cL^{\delta/2}),
\]
which beats every polynomial in \(L\).
The same application with \(\nu_1\) treats the small conductors,
apart from the single exceptional real character.  Abel summation
therefore gives the \(O(L^{-4})\) weighted partial-prime-sum bounds
needed below.  Omitting \(\mathcal L/\nu\) from the intermediate
statement is unnecessary.

\section{Squarefree saddle densities and the covariance correction}
\label{sec:saddle}

For a finite prime set \(J\), write
\[
 \zeta_J^{\rm sf}(s)=\prod_{p\in J}(1+p^{-s}),\qquad
 L_J^{\rm sf}(s,\chi)=\prod_{p\in J}(1+\chi(p)p^{-s}).
\]
In this section \(J=\mathcal J\), with at most \(O(L)\) additional
primes deleted when stated.  The initially omitted primes dividing
\(k_1k_2\) also number \(O(L)\).  Put
\(\phi(s)=\log\zeta_J^{\rm sf}(s)\).  For
\(u_Z=\log Z/\log B\in[u/3,3u]\), define \(\alpha_Z\) by
\[
 -\phi'(\alpha_Z)=\log Z,
 \qquad \lambda_Z=(1-\alpha_Z)\log B,
 \qquad V_Z=\phi''(\alpha_Z).
\]
Monotonicity and the prime number theorem give a unique
\(\alpha_Z\in(4/5,1)\).  Define the central integral
\begin{equation}\label{eq:pure}
 I_{\rm pure}(Z;S)=\frac1{2\pi}
 \int_{-1/\log B}^{1/\log B}
       Z^{\alpha_Z+it-1}\zeta_{J\setminus S}^{\rm sf}(\alpha_Z+it)\,dt.
\end{equation}
This quantity is positive in the range used here.

\begin{lemma}[Uniform saddle expansion]\label{lem:saddle}
Uniformly for \(|S|\le5L\) and the stated \(Z\)-range,
\begin{align}
 \lambda_Z&=\log u_Z+\log\log u_Z+O(1),\label{eq:lambda}\\
 (-1)^j\phi^{(j)}(\alpha_Z)&\asymp_j u_Z(\log B)^j,
               \qquad j\ge1,\label{eq:cumulants}\\
 I_{\rm pure}(Z;S)
 &=\frac{Z^{\alpha_Z-1}\zeta_{J\setminus S}^{\rm sf}(\alpha_Z)}
          {\sqrt{2\pi V_Z}}
       \left(1+\frac{\phi^{(4)}(\alpha_Z)}{8V_Z^2}
                -\frac{5\phi'''(\alpha_Z)^2}{24V_Z^3}
                +O(u_Z^{-2})\right).\label{eq:edgeworth}
\end{align}
Moreover,
\(\left|dI_{\rm pure}(e^x;S)/dx\right|\ll I_{\rm pure}(e^x;S)\).
If \(B^{u/3}\le Z'\le Z\), then
\begin{equation}\label{eq:saddle-ratio}
 \frac{I_{\rm pure}(Z';S')}{I_{\rm pure}(Z;S)}
 \ll (Z/Z')^{1-\alpha_Z}
       \exp\left(O\left(\frac{\log(Z/Z')}{u\log B}\right)\right).
\end{equation}
\end{lemma}

\begin{proof}
Partial summation with input~\ref{input:pnt} gives
\(-\phi'(\alpha)\asymp B^{1-\alpha}/(1-\alpha)\) in this range.
The saddle equation therefore yields
\(e^{\lambda_Z}/\lambda_Z\asymp u_Z\), which gives
\eqref{eq:lambda}; the same calculation with powers of \(\log p\)
gives \eqref{eq:cumulants}.  Deleting \(O(L)\) primes changes the
\(j\)-th derivative by
\(O(LB^{-\alpha_Z/2}(\log B)^j)\), negligible at every scale used.

In \eqref{eq:pure}, put \(z=t\sqrt{V_Z}\).  On
\(|z|\le\log u_Z\), the standardized cumulant of order \(j\) is
\(O_j(u_Z^{1-j/2})\).  Expand the exponential of the cubic through
sixth-order terms.  The odd terms integrate to zero; the fourth term
and the square of the cubic term give the two explicit corrections
in \eqref{eq:edgeworth}.  The remaining Gaussian moments are
\(O(u_Z^{-2})\).  On the rest of the central interval, the elementary
inequality \(1-\cos(t\log p)\gg(t\log p)^2\) gives
\(\exp(-c(\log u_Z)^2)\) relative decay.  This proves the expansion
with a uniform error.

The derivative assertion follows by differentiating the integral
itself, using \(\alpha'(x)=-1/\phi''(\alpha(x))\), and bounding the
resulting \(|t|\) Gaussian moment.  The ratio assertion follows by
integrating the derivative of the explicit logarithmic main term;
\(\alpha'(x)<0\) gives its stated one-sided exponent.
\end{proof}

\subsection{Why the second difference is legitimate}

We will need the first nontrivial covariance between two prime
divisibility indicators.  Let \(x=\log Z\), \(\alpha=\alpha_{e^x}\),
and write
\begin{align*}
 G(x)={}&(\alpha-1)x+\phi(\alpha)
       -\tfrac12\log(2\pi\phi''(\alpha))
       +\log(1+\kappa_1(\alpha)),\\
 \kappa_1(\alpha)={}&
       \frac{\phi^{(4)}(\alpha)}{8\phi''(\alpha)^2}
       -\frac{5\phi'''(\alpha)^2}{24\phi''(\alpha)^3}.
\end{align*}
Then \(\log I_{\rm pure}(e^x;\varnothing)=G(x)+O(u^{-2})\), and
direct differentiation of \(G\) gives
\[
 G''(x)=-\phi''(\alpha)^{-1}(1+O(u^{-1})).
\]
For \(h_1=\log p\), \(h_2=\log p'\), use the exact double-integral
formula for the second difference to obtain
\begin{equation}\label{eq:second-difference}
 G(x-h_1-h_2)-G(x-h_1)-G(x-h_2)+G(x)
       =-\frac{h_1h_2}{V_Z}+O(u^{-2}).
\end{equation}
The four error values from \eqref{eq:edgeworth} contribute only
\(O(u^{-2})\).  No derivative of an unspecified error term is taken.
This distinction is necessary: a uniform bound on a function alone
does not imply a bound on its derivatives.

Let \(\alpha_s=\alpha_{X_{\rm sc}}\) and
\(\mu_i(a)=X_{\rm sc}|P_i(a/A)|\).  Since
\[
 |\log|P_i(a/A)||=O(L^t+d\log d)=O(L^t),
\]
we have
\( |\alpha_{\mu_i(a)}-\alpha_s|\log B=o(u^{-1})\), by \(t<b\).
For a label \(\mathfrak q=(i,p,r_0)\), define
\begin{align*}
 m_{\mathfrak q}(a)
 &=p^{\alpha_s/2-1}
       \frac{I_{\rm pure}(\mu_i(a)/p;\{p\})}
            {I_{\rm pure}(\mu_i(a);\varnothing)},\\
 d_{\mathfrak q}(a)&=p^{\alpha_s/2}m_{\mathfrak q}(a),
 \qquad
 c_{\mathfrak q}(a)=
       \frac{m_{\mathfrak q}(a)\log p}
            {\sqrt{\phi''(\alpha_{\mu_i(a)})}}.
\end{align*}
The expansion gives
\begin{equation}\label{eq:prime-moments}
 m_{\mathfrak q}(a)=p^{-\alpha_s/2}(1+O(u^{-1})),
 \quad d_{\mathfrak q}(a)=1+O(u^{-1}).
\end{equation}
For distinct primes \(p,p'\) in the same field, the exact second-moment
quantity is
\[
 U^{\rm ex}_{\mathfrak q,\mathfrak q'}(a)
 =(pp')^{\alpha_s/2-1}
   \frac{I_{\rm pure}(\mu_i(a)/(pp');\{p,p'\})}
        {I_{\rm pure}(\mu_i(a);\varnothing)}.
\]
Equation~\eqref{eq:second-difference} yields
\begin{equation}\label{eq:covariance}
 U^{\rm ex}_{\mathfrak q,\mathfrak q'}(a)
 =m_{\mathfrak q}(a)m_{\mathfrak q'}(a)
  -c_{\mathfrak q}(a)c_{\mathfrak q'}(a)
  +O((pp')^{-\alpha_s/2}u^{-2}).
\end{equation}
The rank-one negative term is the finite-size constraint imposed by
conditioning a smooth integer to have a specified size.

\section{Clean smooth progressions and the two-prefix estimate}
\label{sec:progressions}

For \(Y,S>0\), put
\[
 W_{Y,S}(x)=S^{-1}\exp(-\pi(x-Y)^2/S^2),\qquad R=Y/S,
\]
and, for \((a_0,q)=1\), define
\begin{equation}\label{eq:sigma}
 \Sigma(Y,S;q,a_0,S_{\rm om})
 =q\sum_{\substack{n\ge1\text{ squarefree}\\
              P(n)\subseteq J\setminus S_{\rm om}\\n\equiv a_0\ (q)}}
              W_{Y,S}(n).
\end{equation}
Here \(P(n)\) denotes the set of prime divisors of \(n\).
When primes dividing \(q\) are removed as well, denote the union of
the removed sets by \(S_{\rm tot}\).

\subsection{A compact cutoff with controlled Mellin tails}

The explicit Gevrey kernel in source Lemmas~8.3--8.4 gives a smooth
cutoff \(\widetilde W_{Y,S}\), supported in
\([e^{-0.31}Y,e^{0.31}Y]\), that agrees with the Gaussian up to an
exponentially small tail.  We use the following precise consequences,
uniformly for \(1/2\le\sigma\le1\): its Mellin transform has
\(L^1\)-norm \(O(RY^{\sigma-1})\); its contribution from
\(|t|>R^2\) is bounded by \(Y^{\sigma-1}\) times a polynomial in
\(R\) times \(\exp(-cR^{0.6})\); and on the central interval its
Mellin transform is
\(Y^{\sigma+it-1}(1+O(R^{-2}))\) after integration against the
symmetric central saddle.  One construction starts with
\(\exp(-1/(1-x^2)^2)\) on \((-1,1)\), extends it by zero, and
normalizes its integral.  The derivative bounds
\(C^{j+1}(j!)^{3/2}\) yield the asserted Fourier decay by repeated
integration by parts.  These source lemmas are parameter-general;
our windows satisfy \(K_1^{0.8}\le R\le K_1^{1.7}\).

Let \(I(Y,S;S_{\rm tot})\) be the central Mellin integral with this
window and Euler product.  The saddle expansion and tail bounds below
give
\begin{equation}\label{eq:I-vs-pure}
 I(Y,S;S_{\rm tot})
 =I_{\rm pure}(Y;S_{\rm tot})(1+O(R^{-2})).
\end{equation}

\subsection{Exceptional conductors}

Let \(\mathcal E_{\rm med}\) consist of primitive characters of
conductor \(d_0\) with
\(D_{\rm small}<d_0\le B^{64}\) having a zero in
\(\Re s\ge1-\nu_0\), \(|\Im s|\le T_*\).
Define \(\mathcal E_{\rm hi}\) similarly for
\(B^{64}<d_0\le Q_{\max}\).  A modulus is called \emph{clean}
if no conductor represented in either set divides it.  Let
\(\mathcal Q_{\rm med}\) be the primes in \(J\) dividing a medium
conductor, and let \(\mathcal P_{\rm gen}^{\rm good}\) be the generator
primes dividing none of those conductors.

These sets occur only in the proof.  The algorithm does not enumerate
characters, zeros, or exceptional conductors.

\begin{lemma}[Exceptional-set size]\label{lem:exception-count}
\begin{align*}
 \log(2+|\mathcal E_{\rm med}|)&=O(L^{b-v})=O(L^{\delta/2}),\\
 \log(2+|\mathcal E_{\rm hi}|)&=O(L^{m_0-v}),\\
 |\mathcal P_{\rm gen}^{\rm good}|
       &=(1/2-o(1))B/\log B.
\end{align*}
Also \(|\mathcal Q_{\rm med}|\le128|\mathcal E_{\rm med}|\).
\end{lemma}

\begin{proof}
Apply \eqref{eq:density} with \(\sigma=1-\nu_0\) and
\(T=T_*\).  The logarithm of \(Q^2T_*\) is \(O(L^b)\) at
\(Q=B^{64}\), and \(O(L^{m_0})\) at \(Q=Q_{\max}\).
The logarithmic factor in \eqref{eq:density} is smaller than both
displayed positive powers.  Each medium conductor has at most
128 prime divisors greater than \(\sqrt B\).  Removing their union
therefore deletes \(o(B/\log B)\) primes from \((B/2,B]\).
\end{proof}

\subsection{Euler-product decay}

For the principal character at \(|t|\ge1/\log B\), and for the
exceptional real character at every \(t\), the required estimate is
\begin{equation}\label{eq:repulsion}
 \frac{|L_{J\setminus S}^{\rm sf}(\alpha_Y+it,\chi)|}
      {\zeta_{J\setminus S}^{\rm sf}(\alpha_Y)}
 \ll\exp\left(-\frac{cu_Y}{(\log u_Y)^2}\right),
 \qquad |t|\le R^2.
\end{equation}
Here \(|S|\le5L\); the principal-character restriction on \(t\)
is necessary.

To see the low-frequency estimate explicitly, the logarithm of the
Euler-product ratio is at most
\[
 -c\sum_{\sqrt B<p\le B}p^{-\alpha_Y}
       (1-\Re(\chi(p)p^{-it}))+O(LB^{-\alpha_Y/2}).
\]
The explicit formula from input~\ref{input:page} gives, on
\([\sqrt B,B]\), a prime-counting main term with derivative \(1\)
for the principal character and derivative
\(-\mathbf1_{\{\beta\}}x^{\beta-1}\le0\) for the exceptional
character.  At \(|t|\le L^{50}/\log B\), partial summation changes
the displayed sum into an integral in \(v=\log x/\log B\), with
weight \(e^{\lambda_Yv}/v\) and negligible error.
On \([1-1/\lambda_Y,1]\), its oscillatory factor is bounded below
by a constant multiple of \(\sin^2(t\log B\,v/2)\) in the principal
case or \(\cos^2(t\log B\,v/2)\) in the exceptional case.
Integrating over that interval gives at least
\(c e^{\lambda_Y}/\lambda_Y^3\asymp cu_Y/\lambda_Y^2\).
This covers short intervals about zeros of the sine or cosine as well
as intervals containing whole periods.

For \(L^{50}/\log B<|t|\le R^2\), the zero-free strip and
Lemma~\ref{lem:prime-sum} bound the weighted character prime sum by
\(O(L^{-4})\).  The untwisted prime mass is \(\asymp u_Y\), giving
the stronger decay \(\exp(-cu_Y)\).  An exceptional real zero lies
outside these shifted rectangles.  This proves
\eqref{eq:repulsion}, including a possible exceptional character
without a lower bound on \(1-\beta\).

\begin{proposition}[Clean progression asymptotic]\label{prop:progressions}
Uniformly for clean \(1\le q\le Q_{\max}\), \((a_0,q)=1\),
\(|S_{\rm om}|\le2L\), \(K_1^{0.8}\le Y/S\le K_1^{1.7}\), and
\(u_Y\in[u/2,2u]\),
\begin{equation}\label{eq:clean-asymptotic}
 \Sigma(Y,S;q,a_0,S_{\rm om})
 =\frac q{\varphi(q)}I(Y,S;S_{\rm tot})
       \left(1+O(\exp(-c_\delta L^\eta))\right).
\end{equation}
\end{proposition}

\begin{proof}
Character orthogonality gives the principal-character contribution
\(qI/\varphi(q)\).  The cutoff and \eqref{eq:repulsion} control
the exceptional real character, if it induces a character modulo
\(q\).  There is at most one such induced character.

The other characters are handled by two prime prefixes.  Since
\begin{equation}\label{eq:prefix-size}
 q^2B^2<e^{-0.31}Y,
\end{equation}
every squarefree smooth integer in the cutoff support has a unique
ordered factorization \(n=m_1m_2m_3\), in which \(m_1\) is the
shortest initial prime prefix exceeding \(q\), \(m_2\) is the next
shortest such prefix, and
\(q<m_1,m_2\le qB\).  Let \(p_a<p_b\) be their final primes.
For fixed \(p_a,p_b\), let \(A_{p_a}(s,\chi)\) and
\(B_{p_a,p_b}(s,\chi)\) be the two prefix Dirichlet polynomials.
The remaining factor contributes the squarefree Euler product over
primes above \(p_b\).

Shortestness gives \(m_1/p_a\le q\); distinct possible quotients
are distinct reduced residues modulo \(q\).  Consequently
\[
 \sum_{\chi\bmod q}|A_{p_a}(\alpha+it,\chi)|^2
 =\varphi(q)\sum_{m_1}m_1^{-2\alpha}
 \le\varphi(q)^2q^{-2\alpha},
\]
and the same bound holds for the second prefix.  Cauchy--Schwarz
therefore gives the decisive estimate
\begin{equation}\label{eq:prefix-cs}
 \frac q{\varphi(q)}
       \sum_{\chi\bmod q}|A_{p_a}B_{p_a,p_b}|
 \le q^{2(1-\alpha)}.
\end{equation}

For a remaining character, its primitive conductor avoids the
excluded zero sets.  The corrected prime-sum estimate gives
\(\exp(-cu_Y)\) relative decay in the terminal Euler product if
\(p_b\le B^{1-1/(4\lambda_Y)}\).  If \(p_b\) exceeds this threshold,
the omitted lower primes already carry a fixed positive fraction
of the saddle prime mass, giving the same bound.  There are at
most \(B^2\) terminal-prime pairs.  The cutoff's Mellin \(L^1\)
bound and the saddle main term show that the nonexceptional relative
error is at most
\begin{equation}\label{eq:progression-error}
 CB^2R\sqrt{V_Y}
       \exp(2(1-\alpha_Y)\log q-cu_Y).
\end{equation}
The exceptional-character relative error is at most
\(CR\sqrt{V_Y}\exp(-cu_Y/(\log u_Y)^2)\).
The Mellin tail is negligible because its
\(\exp(-cR^{0.6})\) factor dominates all losses exponential in
\(\log q+u_Y\).

The positive logarithms in \eqref{eq:progression-error} total
\[
 O(L^b+L^k+L^{m_0-b}L_2+\log L)=o(L^r)=o(u_Y).
\]
Also \(u_Y/(\log u_Y)^2\gg L^\eta\gg\log R\), and
\eqref{eq:prefix-size} holds since
\(2\log q+2\log B=O(L^{m_0})=o(\log Y)\).
These prove \eqref{eq:clean-asymptotic} and
\eqref{eq:I-vs-pure} with the stated uniformity.
\end{proof}

\section{Conditional Gaussian laws at one and two evaluation points}
\label{sec:gaussian}

For each field, give \(\mathbf u=(u_0,u_1)\in\mathbb Z^2\) probability
proportional to
\[
 \exp\left(-\pi\frac{(u_0-\bar u_0)^2}{S_0^2}
           -\pi\frac{(u_1-\bar u_1)^2}{S_1^2}\right).
\]
Call this distribution \(\mathbb P_G\), and use independent copies for
the two fields.  Its mass outside \(\mathcal B_U\) is at most
\(4\exp(-\pi L_0^{40})\).

\subsection{An upper bound that requires no cleanliness}

\begin{lemma}[Greedy-prefix Rankin bound]\label{lem:rankin}
Suppose \((a_0,q)=1\), \(Y/S>1\), \(Y>q\), and
\(S>e^{0.62}qB\).  For \(0<\alpha\le1\),
\[
 \Sigma(Y,S;q,a_0,S_{\rm om})
 \ll (qB/S)^{1-\alpha}\zeta_{J\setminus S_{\rm om}}^{\rm sf}(\alpha)
       +\exp(-c(Y/S)^2).
\]
The same bound holds with primes dividing \(q\) removed from the
Euler product.
\end{lemma}

\begin{proof}
Outside \([e^{-0.31}Y,e^{0.31}Y]\), the normalized Gaussian mass on a
progression is exponentially small.  Within that interval take the
shortest prime prefix \(d_1\) of a supported integer \(n\) that is at
least \(M_1=e^{0.31}Yq/S\).  Then \(d_1<M_1B\), and its complement
\(d_2=n/d_1\) lies in
\[
 e^{-0.62}S/(qB)<d_2\le S/q.
\]
For fixed \(d_2\), \(d_1\) occupies one class modulo \(q\), and its
normalized Gaussian mass is at most \(3/d_2\).  Since
\[
 d_2^{-1}\le(e^{-0.62}S/(qB))^{\alpha-1}d_2^{-\alpha},
\]
summing over squarefree smooth \(d_2\) proves the assertion.
\end{proof}

The cancellation of \(Y\) in this bound is useful.  Dividing a forced
gcd out of a progression does not introduce an extra unaccounted factor
\(G^{1-\alpha}\): the ratio of width to modulus is unchanged.

\subsection{Conditioning on prescribed roots}

Fix a field \(i\) and a set
\(R_i=\{(p_j,r_j)\}_{j=1}^{n_i}\), where \(n_i\le2\), the primes
are distinct members of \(J\setminus\{\ell\}\), and
\(r_j\ne m\pmod{p_j}\).  Let \(\mathcal E_{R_i}\) be the event
that all these residues are roots of the perturbed polynomial.
Each root condition is precisely
\begin{equation}\label{eq:root-congruence}
 u_0+r_ju_1\equiv c_j\pmod{p_j},
\end{equation}
because \(p_j\nmid P_0\ell(r_j-m)\).
Write
\[
 Q_R=\prod_{j=1}^{n_i}p_j,\quad
 Q_{\rm hit}(a)=\prod_{a\equiv r_j\ (p_j)}p_j,
 \quad Q_{\rm pre}(i,a)=\ell^{\epsilon_\ell(i,a)}Q_{\rm hit}(a),
\]
where \(\epsilon_\ell(i,a)=1\) exactly when \(i=1\) and \(a\)
is one of the prescribed roots modulo \(\ell\).
Let \(E_i(a)\) denote the event that \(|f_i(a)|\) is a nonzero,
squarefree \(J\)-smooth integer.

Put
\[
 L_a(\mathbf u)=u_0+au_1,\quad
 \bar\ell_a=\bar u_0+a\bar u_1,\quad
 S(a)=\sqrt{S_0^2+a^2S_1^2},\quad
 \widetilde S_1(a)=\frac{S_0S_1}{S(a)}.
\]
For \(a\in\mathcal A\), \(S(a)\asymp UA\) and
\(\widetilde S_1(a)\asymp U\).
The change of variables \((u_0,u_1)\mapsto(L_a,u_1)\) is
unimodular.  Completing the square gives the exact identity
\begin{align}\label{eq:gaussian-square}
 \frac{(u_0-\bar u_0)^2}{S_0^2}
 +\frac{(u_1-\bar u_1)^2}{S_1^2}
 ={}&\frac{(L_a-\bar\ell_a)^2}{S(a)^2}\notag\\
 &+\frac{\left(u_1-\bar u_1-
      \frac{aS_1^2}{S(a)^2}(L_a-\bar\ell_a)\right)^2}
        {\widetilde S_1(a)^2}.
\end{align}
Poisson summation on the remaining coset proves
\begin{equation}\label{eq:root-probability}
 \mathbb P_G(\mathcal E_{R_i})=Q_R^{-1}(1+O(e^{-cU})),
\end{equation}
and, conditionally, \(L_a\) has a discrete Gaussian on one class modulo
\(Q_{\rm hit}(a)\), with mean \(\bar\ell_a\), width \(S(a)\), and
normalizing factor \(Q_{\rm hit}(a)/S(a)\), up to relative
\(O(e^{-cU})\).  Indeed, every residual width divided by its modulus
is at least \(U/(\sqrt2B^2)\), and \(U\ge4B^4\) eventually.

Let \(s_i(a)=\operatorname{sgn}P_i(a/A)\), and write
\[
 \mu_i^{\rm ex}(a)=s_i(a)
   \bigl(f_i^{(0)}(a)+P_0\ell(a-m)\bar\ell_a\bigr),\qquad
 S_{1,i}^*(a)=P_0\ell|a-m|S(a).
\]
The construction bounds give
\(\mu_i^{\rm ex}(a)=\mu_i(a)(1+O(K_1^{-0.95}))\).
On the overwhelmingly likely central range of \(L_a\), the normalized
value \(s_i(a)f_i(a)/Q_{\rm pre}(i,a)\) runs through one residue
class modulo
\begin{equation}\label{eq:q1}
 q_1=P_0\ell^{1-\epsilon_\ell(i,a)}|a-m|,
\end{equation}
with center \(\mu_i^{\rm ex}(a)/Q_{\rm pre}(i,a)\) and width
\(S_{1,i}^*(a)/Q_{\rm pre}(i,a)\).  Its width/modulus ratio is
\(S(a)/Q_{\rm hit}(a)\ge UA/B^2\).

The class in \eqref{eq:q1} is coprime.  The factors of \(P_0\) are
excluded by \(f_i(a)\equiv1\pmod{P_0}\).  A factor of \(a-m\)
cannot divide \(f_i(a)\), because
\(f_i(a)\equiv k_iN\pmod{a-m}\) and
\((a-m,k_iN)=1\).  The \(\ell\)-case is exactly the distinction in
\(\epsilon_\ell\).  Also, prescribed roots are different from
\(m\), so dividing their prime factors does not spoil this argument.

\subsection{The two-point lattice, including repeated prime factors}

For \(a'\ne a\), put \(\Delta=|a'-a|\).  The map
\(\mathbf u\mapsto(L_a,L_{a'})\) identifies \(\mathbb Z^2\) with
\[
 \{(x,y)\in\mathbb Z^2:y\equiv x\pmod\Delta\}.
\]
This lattice relation prevents one from treating the two values as
independent.  Put
\[
 Q_{\rm both}=(Q_{\rm hit}(a),Q_{\rm hit}(a')),
 \quad Q_{\rm free}=Q_R/Q_{\rm hit}(a),
 \quad Q_{\rm free}^*=Q_{\rm free}/(Q_{\rm free},Q_{\rm hit}(a')).
\]
Given \(L_a\), the exact step for \(L_{a'}\) is
\begin{equation}\label{eq:M2}
 M_2=\frac{\Delta}{(\Delta,Q_{\rm hit}(a'))}
          Q_{\rm hit}(a')Q_{\rm free}^*.
\end{equation}
This formula follows prime by prime from \eqref{eq:root-congruence}.
If a prescribed prime hits both points, it divides \(\Delta\), and
its existing power in \(\Delta\) is retained rather than duplicated.
If it hits neither, the root condition adds one power to the
\(\Delta\)-step.  If it hits just one, the two points are distinct
modulo that prime and the one relevant congruence is imposed.

The conditional mean of \(L_{a'}\), given \(L_a=x\), is
\[
 \bar\ell_{a'}+\frac{S_0^2+aa'S_1^2}{S(a)^2}(x-\bar\ell_a),
\]
and its width is \(\Delta\widetilde S_1(a)\).
After normalizing the polynomial value by its prescribed factors,
the second modulus is
\begin{equation}\label{eq:q2}
 q_2=P_0\ell^{1-\epsilon_\ell(i,a')}|a'-m|
             \frac{\Delta}{Q_{\rm both}}Q_{\rm free}^*.
\end{equation}
The width/modulus ratio is at least \(U/(\sqrt2B^2)\).
The conditional center differs relatively from
\(\mu_i(a')/Q_{\rm pre}(i,a')\) by at most \(K_1^{-0.94}\),
when \(|L_a-\bar\ell_a|\le L_0^3S(a)\).

For separated points \(\Delta>\Delta_{\rm close}\), the first and
second center/width ratios lie in
\([K_1^{0.8},K_1^{1.7}]\).  These assertions follow by substituting
\eqref{eq:scale}: the ratios are \(K_1^{1+o(1)}\) and, at worst,
\(K_1^{3/2+o(1)}\).  The factors \(|P_i(a/A)|\), polynomial powers
of \(L\), and coefficient errors are \(K_1^{o(1)}\), since \(t<k\).
All normalized centers remain in the saddle range.

\subsection{Cleanliness and exact local factors}

Call \(a\) \emph{one-clean} for \((i,R_i)\) when
\(P_0\ell|a-m|Q_R\) is a clean modulus.  Call \((a,a')\)
\emph{two-clean} when
\begin{gather*}
 \Delta>\Delta_{\rm close},\qquad \ell^2\nmid\Delta,
 \qquad p_j^2\nmid\Delta\quad((p_j,r_j)\in R_i),\\
 P_0\ell|a-m||a'-m|\Delta Q_R
       \text{ is clean}.
\end{gather*}
The product in the second line is at most \(Q_{\max}\) eventually:
its logarithm is \(2L^{m_0}+o(L^{m_0})\).
Both \eqref{eq:q1} and \eqref{eq:q2}, and the enlarged moduli used
to exclude additional bad residues, divide this product.

Define the reference densities
\begin{align*}
 q_0(a)&=P_0\ell|a-m|,\qquad D_0(a)=J\cap\{p:p\mid a-m\},\\
 Z_i(a)&=\frac{q_0(a)}{\varphi(q_0(a))}
     I(\mu_i(a),S_{1,i}^*(a);\{\ell\}\cup D_0(a)),\\
 \Phi_i(a;R_i)&=
    \frac{I_{\rm pure}(\mu_i(a)/Q_{\rm pre}(i,a);
                          P(Q_{\rm pre}(i,a)))}
         {I_{\rm pure}(\mu_i(a);\varnothing)}.
\end{align*}
Put
\begin{equation}\label{eq:gamma}
 \gamma(a,a')=
 \prod_{\substack{z_0<p\le\sqrt B\\p\mid a'-a,\ p\nmid a-m}}
                \frac p{p-1}.
\end{equation}

\begin{proposition}[Conditional smoothness laws]\label{prop:conditional}
There is an absolute \(C\) such that, with
\(R_{\rm rough}=\exp(CL^{m_0-b}L_2)\), the conditional one-point
probability is at most
\(R_{\rm rough}Z_i(a)\Phi_i(a;R_i)\), and the conditional
two-point probability is at most
\(R_{\rm rough}^2Z_i(a)Z_i(a')\Phi_i(a;R_i)\Phi_i(a';R_i)\).
For one-clean points,
\[
 \mathbb P_G(E_i(a)\mid\mathcal E_{R_i})
   =Z_i(a)\Phi_i(a;R_i)(1+O(K_1^{-0.90})).
\]
For two-clean pairs,
\begin{align*}
 &\mathbb P_G(E_i(a)\cap E_i(a')\mid\mathcal E_{R_i})\\
 &\quad=Z_i(a)Z_i(a')\Phi_i(a;R_i)\Phi_i(a';R_i)
          \gamma(a,a')(1+O(K_1^{-0.90})).
\end{align*}
\end{proposition}

\begin{proof}
For a single point, apply Lemma~\ref{lem:rankin} or
Proposition~\ref{prop:progressions} to \eqref{eq:q1}.  For the second
point condition first on \(L_a\), then use \eqref{eq:q2}.
If its residue has gcd \(G\) with \(q_2\), a squarefree smooth event
forces \(G\) itself to be squarefree and \(J\)-smooth.  Otherwise
the event is empty.  Divide \(G\) out and exclude its primes from
the remaining value.  The width/modulus ratio is unchanged, so
Lemma~\ref{lem:rankin} still applies.  Here \(G\le AB^3\), and
the more precise saddle comparison costs at most
\[
 \exp\left(O\left(\frac{(\log G)^2}{u(\log B)^2}\right)\right)
 =1+o(1),
\]
because \(2c<h+b\).
The logarithm of the remaining rough loss is bounded by
\[
 O\left((1-\alpha)\log\frac{\mu_i(a)}{UA}\right)
       =O(L^{m_0-b}L_2).
\]
The \(U\)-contribution cancels in this ratio by \eqref{eq:scale};
using \(\log\mu_i(a)\) alone would give the wrong error budget.

For the clean two-point formula, inspect each prime dividing \(q_2\).
Primes dividing \(P_0|a'-m|\) cannot divide its normalized residue.
The prime \(\ell\) can introduce one additional forbidden residue
when both points hit its root classes; the condition
\(\ell^2\nmid\Delta\) makes the two squarefreeness exclusions
distinct.  Its probability is \(O(\ell^{-1})\) times the first-point
main term.  Primes in \((z_0,\sqrt B]\) dividing \(\Delta\) and
not \(a-m\) contribute exactly \(p/(p-1)\), giving
\eqref{eq:gamma}.  A factor-base prime dividing \(\Delta\) removes
only \(O(p^{-\alpha})\) of the first-point smooth mass.  Each of
the at most two other prescribed primes likewise removes
\(O(p^{-\alpha})\); for a zero residue divide that prime out, and
for a nonzero residue enlarge the modulus.  All these moduli divide
the two-clean product.

The total discarded relative mass is at most
\[
 L^{O(1)}\exp(O(L^{m_0-b}L_2))B^{-\alpha/2}
       =o(K_1^{-0.90}),
\]
since \(b>m_0-b\) and \(b>k\).
The center perturbations cost \(O(K_1^{-0.94})\) by direct
differentiation of the saddle integral.  Gaussian tails and Poisson
errors are smaller still.  Integrating over \(L_a\) proves both
clean formulas.
\end{proof}

\section{Centered Gram moments and exceptional-conductor geometry}
\label{sec:moments}

Fix \(\ell\in\mathcal P_{\rm gen}^{\rm good}\).  A factor-base label
is a triple \(\mathfrak q=(i,p,r_0)\), where \(i\in\{1,2\}\),
\(p\in J\), and \(r_0\) is a root of \(f_i\) modulo \(p\).
Let \(F_{\rm all}\) be the set of these labels.  Let
\(F_{\rm mult}\) be the labels for multiple roots, and let
\(F_{\rm shared}\) be those at a residue that is a root of both
polynomials modulo the same prime.  Set
\[
 F_{\rm clean}=F_{\rm all}\setminus(F_{\rm mult}\cup F_{\rm shared}),
 \quad
 F_{\rm med}=\{(i,p,r_0)\in F_{\rm all}:p\in\mathcal Q_{\rm med}\},
 \quad F^\circ=F_{\rm all}\setminus F_{\rm med}.
\]
The \(d\) protected labels are
\(\mathfrak l_s=(1,\ell,r_s)\), \(1\le s\le d\).  They belong
to \(F^\circ\cap F_{\rm clean}\) for every perturbation.

For a set \(D\) of labels, let \(\mathcal S(D)\) consist of those
\(a\in\mathcal A^*\) for which both absolute polynomial values are
squarefree \(J\)-smooth and which hit no label of \(D\).
Write \(\mathcal S_{\rm raw}=\mathcal S(\varnothing)\).
For \(a\in I_j\cap\mathcal A^*\), put
\begin{equation}\label{eq:weights}
 \pi_0(a)=\frac1{(2d+1)|I_j\cap\mathcal A^*|},\qquad
 w_{\rm eq}(a)=\frac{\pi_0(a)}{Z_1(a)Z_2(a)}.
\end{equation}
In particular \(\sum_a\pi_0(a)=1\).
Let
\[
 z_{\mathfrak q}(a)=p^{\alpha_s/2}\mathbf1_{a\equiv r_0\ (p)},
 \quad
 \boldsymbol\ell_i(a)=
       (\log|a/A-\rho_{i,s}|)_{s=1}^d,
 \quad \boldsymbol\ell=(\boldsymbol\ell_1,\boldsymbol\ell_2).
\]
The unprojected feature vector is
\[
 \mathbf r^\circ(a)=
   (1,(z_{\mathfrak q}(a))_{\mathfrak q\in F^\circ},
                  \boldsymbol\ell(a)),
\]
and the random Gram matrix and random total weight are
\[
 \widetilde H^\circ=
      \sum_{a\in\mathcal S_{\rm raw}}
           w_{\rm eq}(a)\mathbf r^\circ(a)\mathbf r^\circ(a)^T,
 \qquad
 \widehat C=\sum_{a\in\mathcal S_{\rm raw}}w_{\rm eq}(a).
\]

\subsection{A fixed index set for a random factor base}

Moment conditioning must be meaningful before the polynomials are
sampled.  Use the universal label set consisting of the protected
labels and every \((i,p,r_0)\) with
\(p\in J\setminus(\mathcal Q_{\rm med}\cup\{\ell\})\) and
\(r_0\ne m\pmod p\).  A label enters \(F^\circ\) precisely when
its root event occurs.  We condition on that event when estimating
an entry.  The factor \(1/p\) from
\eqref{eq:root-probability} is retained when summing over universal
indices.  This avoids paying the cardinality of that universal set
as a spurious matrix-dimension factor.

On the universal index set, start with the rank-one matrix whose mean
vector is
\((1,(m_{\mathfrak q}(a)),\boldsymbol\ell(a))\).
Replace each prime diagonal by \(d_{\mathfrak q}(a)\), and replace
same-field, distinct-prime entries by the exact second-moment quantity
in \eqref{eq:covariance}.  Call the resulting matrix
\(U^{\rm ex}(a)\), and set
\[
 U^{\rm ex}=\sum_a\pi_0(a)U^{\rm ex}(a),\qquad
 \widetilde A^\circ=
       \widetilde H^\circ-\widehat C\,U^{\rm ex}|_{F^\circ}.
\]
The restriction includes the scalar and archimedean rows.  Centering
by \(\widehat C\) is essential: it removes the random total relation
mass before the entrywise cancellation argument.

Define size weights
\[
 M_0=1,\quad M_{\mathfrak q}=p^{-\alpha_s/2},\quad
 M_\nu=L_0\quad(\nu\text{ archimedean}),
\]
and \(M_{\alpha\beta}=M_\alpha M_\beta\), except that
\(M_{\mathfrak q\mathfrak q}=1\).

\begin{proposition}[Centered moments]\label{prop:moments}
For all sufficiently large \(N\),
\[
 \mathbb E\widehat C=1+O(K_1^{-0.90}),\qquad
 \operatorname{Var}(\widehat C)=O(L^{-a}).
\]
For each pair of indices \((\alpha,\beta)\), except distinct labels
with the same underlying rational prime, let
\(\mathcal E_{\alpha\beta}\) require their root events.  Then
\begin{equation}\label{eq:centered-moment}
 \mathbb E\left[(\widetilde A^\circ_{\alpha\beta})^2
             \mid\mathcal E_{\alpha\beta}\right]
 \le M_{\alpha\beta}^2E.
\end{equation}
\end{proposition}

\begin{proof}
There are at most two prescribed roots in each field.  Let \(R_i\)
record those roots, put
\(\Pi_R(a)=\Phi_1(a;R_1)\Phi_2(a;R_2)\), and define
\[
 f_*(a)=r^\circ_\alpha(a)r^\circ_\beta(a),\qquad
 g_*(a)=f_*(a)-U^{\rm ex}_{\alpha\beta}.
\]
The centered square moment is exactly
\begin{equation}\label{eq:moment-sum}
 \sum_{a,a'\in\mathcal A^*}
   w_{\rm eq}(a)w_{\rm eq}(a')g_*(a)g_*(a')
   \prod_{i=1}^2
      \mathbb P_G(E_i(a)\cap E_i(a')\mid\mathcal E_{R_i}).
\end{equation}
Partition the sample points into \(O(d)\) cells according to the
prescribed prime and protected-root classes they hit.  On each cell,
the hit-prime product \(M_c\le B^3\) fixes a single residue class,
and
\begin{equation}\label{eq:cell-bounds}
 |g_*(a)|\Pi_R(a)\ll L^{O(1)}M_cM_{\alpha\beta},\qquad
 \sum_a\pi_0(a)|g_*(a)|\Pi_R(a)\ll L^{O(1)}M_{\alpha\beta}.
\end{equation}
This follows directly from the definitions and
Lemma~\ref{lem:brun}; a hit contributes
\(p^{\alpha_s/2}p^{1-\alpha_s}=p\,p^{-\alpha_s/2}\), or \(p\)
for a prime diagonal.

We first bound the exceptional pairs in \eqref{eq:moment-sum}.
The diagonal \(a=a'\) contributes an exponentially small quantity
on the \(L^c\)-scale: its inverse interval size supplies
\((Aw)^{-1}\), while
\[
 \log((Z_1(a)Z_2(a))^{-1})=O(u\log u),
\]
and all remaining losses have logarithms
\(O(L^b+L^k+L^t+L^{m_0-b}L_2)\).  Here \(c>r,b,k,t\).
Close pairs \(0<|a'-a|\le\Delta_{\rm close}\) contribute
\(K_1^{-1/2+o(1)}M_{\alpha\beta}^2\).  If a prescribed prime or
\(\ell\) has its square dividing \(a'-a\), one extra factor of
that prime occurs in the progression modulus, giving
\(B^{-1/2+o(1)}M_{\alpha\beta}^2\).

For clarity, the conductor calculations are given in detail.  The
elementary bound
\begin{equation}\label{eq:divisor-weight}
 \sum_{d\mid n}d^{-1/2}\le\exp(C\sqrt{\log n})
\end{equation}
follows by splitting the Euler product at \(\log n\) and summing
\(p^{-1/2}\) below that point.  It is a weighted divisor bound;
using the full divisor count would spend an unnecessary loss.

For a medium conductor dividing
\(P_0\ell|a-m||a'-m||a'-a|Q_{\rm pair}\), its part outside
\(P_0\) has size at least \(D_{\rm small}^{0.99}\).  It is
coprime to the prescribed and protected primes, by the definitions
of \(F^\circ\) and \(\mathcal P_{\rm gen}^{\rm good}\).
One of the three differences therefore has a divisor at least
\(D_{\rm small}^{0.33}\).  Take its shortest prime prefix above
that threshold.  If this prefix is at most \(B^{10}\), progression
counting saves
\(D_{\rm small}^{-0.165}\) times the inverse square root of the
prefix.  If it exceeds \(B^{10}\), its last prime exceeds \(B^9\),
giving an even stronger saving.  Count in \(a\) or in \(a'\),
as appropriate to the difference being divided.  After
\eqref{eq:divisor-weight}, the logarithmic costs are
\[
 O(L^{b-v}+L^{b/2}+L^{m_0-b}L_2),
\]
whereas the saving is \(cL^b/L_2\).  Each cost is smaller by a
strict power gap.

For a high conductor, removing \(P_0\ell Q_{\rm pair}\) leaves a
factor larger than \(B^{60}\).  One difference has a divisor above
\(B^{20}\).  Its shortest prefix above \(B^{10}\) either lies
below \(B^{20}\), saving
\(B^{-2}\) times the inverse square root of that prefix after
allowing a gcd of size \(B^3\), or has a last prime above
\(B^{10}\), giving a stronger saving.  The total logarithmic cost
is
\[
 O(L^{m_0-v}+L^{m_0/2}+L^{m_0-b}L_2)=o(L^b).
\]
This proves a fixed negative power of \(B\) for the high-conductor
contribution.  For prefixes up to \(B^{10}\) or \(B^{20}\), plain
progression counting suffices: \(Aw\) dominates every fixed power
of \(B\).  Thus Lemma~\ref{lem:brun} is not being applied outside
its stated \(qQ\le B^4\) range.

On the remaining two-clean pairs, use
Proposition~\ref{prop:conditional}.  Apart from negligible errors,
\eqref{eq:moment-sum} becomes
\begin{equation}\label{eq:gamma-main}
 \sum_{a,a'}\pi_0(a)\pi_0(a')g_*(a)g_*(a')
                      \Pi_R(a)\Pi_R(a')\gamma(a,a')^2.
\end{equation}
Expand the last factor as
\[
 \gamma(a,a')^2
 =\sum_h\eta(h)\mathbf1_{h\mid a'-a,\ (a-m,h)=1},
 \quad
 \eta(h)=\prod_{p\mid h}\frac{2p-1}{(p-1)^2},
\]
over squarefree \(h\) supported on \((z_0,\sqrt B]\).
The tail \(h>\sqrt B\), treated by its shortest prefix in
\((\sqrt B,B]\), costs \(B^{-0.49}M_{\alpha\beta}^2\).

For \(h\le\sqrt B\), replace the protected-prime hit in
\(\Pi_R\) by the one forced by \(f_*\), obtaining
\(\Pi_R^\circ\).  The difference has relative weighted mass
\(O(L^{O(1)}B^{-0.49})\).  Let \(Q_{\rm pair}\) be the product
of distinct rational primes represented by the indices
\(\alpha,\beta\), and let \(r_{\rm pair}\) be their simultaneous
residue.  The crucial exact identity is
\begin{equation}\label{eq:exact-cancellation}
 f_*(a)\Pi_R^\circ(a)
   =\mathbf1_{a\equiv r_{\rm pair}\ (Q_{\rm pair})}
             Q_{\rm pair}U^{\rm ex}_{\alpha\beta}(a).
\end{equation}
It holds for scalar and archimedean indices, one-prime indices,
prime diagonals, distinct primes in one field, and primes in
different fields.  For example, on a prime diagonal it is the
identity \(p^{\alpha_s}\Phi=p\,d_{\mathfrak q}\); for two
distinct primes it is exactly the definition of the second-moment
ratio.  Replacing \(U^{\rm ex}\) by a heuristic product of means
would destroy this identity.

Apply Lemma~\ref{lem:brun} on a fixed reduced class modulo \(h\)
to \eqref{eq:exact-cancellation}, and separately to
\(\Pi_R^\circ\).  The main terms are respectively
\(U^{\rm ex}_{\alpha\beta}/\psi_0(h)\) and
\(1/\psi_0(h)\), so they cancel in the sum with \(g_*\).
The functions involved have the required Lipschitz bounds:
differentiate their defining saddle integrals and explicit means,
rather than an asymptotic remainder.  The resulting polynomial
factor in \(d\) is well below the allowed \(L^{40}\).
Thus each class contributes at most
\(L^{O(1)}M_{\alpha\beta}E_{\rm Br}/\psi_0(h)\).
Since \(\varphi(h)\le\psi_0(h)\) and
\(\sum_h\eta(h)/\psi_0(h)\ll1\),
\eqref{eq:gamma-main} is
\(O(L^{O(1)}M_{\alpha\beta}^2E_{\rm Br}^2)\).
By \eqref{eq:errors}, this is \(o(M_{\alpha\beta}^2E)\).

The same one-point decomposition proves the expectation of
\(\widehat C\).  With no feature indices, the expansion of
\(\gamma^2-1\) gives variance
\(O(\sum_{p>z_0}p^{-2}+K_1^{-0.48})=O(L^{-a})\).
This proves the proposition.
\end{proof}

\section{Adaptive deletion with a deterministic invariant}
\label{sec:peeling}

A small average entrywise error does not automatically give a useful
Gram matrix: a few labels may carry too much error or appear only in
relations that must be removed.  We delete such labels and then delete
any label whose surviving mass has fallen too far.  The argument must
control this adaptive process without a union bound over its stages.

For a label set \(D\), write
\begin{align*}
 \mu(D)&=\sum_{\mathfrak q=(i,p,r_0)\in D\cap F^\circ}p^{-\alpha_s},\\
 \Delta_{\mathfrak q}(D)&=
   \sum_{a\in\mathcal S_{\rm raw}\setminus\mathcal S(D)}
             w_{\rm eq}(a)z_{\mathfrak q}(a)^2,\\
 W_{\rm drop}(D)&=
   \sum_{a\in\mathcal S_{\rm raw}\setminus\mathcal S(D)}w_{\rm eq}(a).
\end{align*}
Choose
\begin{equation}\label{eq:peel-thresholds}
 \tau=E^{1/4},\qquad \zeta=E^{1/8},\qquad \mu_*=E^{7/15}.
\end{equation}
The symbol \(\zeta\) in this section is a deletion threshold, not
an Euler product.  Let
\[
 F_{\rm diag}=\{\mathfrak q\in F^\circ:
                  |\widetilde A^\circ_{\mathfrak q\mathfrak q}|>\tau\},
 \quad
 D_0=F_{\rm med}\cup F_{\rm mult}\cup F_{\rm shared}\cup F_{\rm diag},
\]
and recursively set
\begin{equation}\label{eq:peel-process}
 D_{j+1}=D_j\cup
       \{\mathfrak q\in F^\circ\setminus D_j:
                               \Delta_{\mathfrak q}(D_j)>\zeta\}.
\end{equation}
There are finitely many labels, so this process terminates.

\subsection{The good event and its probability}

Let \(F_A^2\) be the sum of squared centered entries, omitting prime
diagonals and pairs with the same underlying prime.  Let
\[
 D_A=\sum_{\mathfrak q\in F^\circ}
           p^{-\alpha_s}(\widetilde A^\circ_{\mathfrak q\mathfrak q})^2,
 \qquad
 F_\ell^2=\sum_{s=1}^d
      \sum_{\substack{\mathfrak q\in F^\circ\\p\ne\ell}}
          (\widetilde A^\circ_{\mathfrak l_s,\mathfrak q})^2.
\]
Conditioned root probabilities give
\begin{align}
 \mathbb E D_A&\ll uE,&
 \mathbb E F_A^2&\ll L^{O(1)}E,&
 \mathbb E F_\ell^2&\ll du\ell^{-\alpha_s}E.\label{eq:aggregate-errors}
\end{align}
For example, summing over all \(r_0\pmod p\) supplies \(p\)
possible labels, but each has root probability \(1/p+O(e^{-cU})\).
The net weight is \(p^{-\alpha_s}\), whose prime sum is \(O(u)\).
This is why no factor of the ambient dimension \(B\) appears in
\eqref{eq:aggregate-errors}.

The two congruences imposing a multiple root have probability
\(p^{-2}(1+O(e^{-cU}))\).  The same holds for a shared root, by
independence of the fields.  Their expected \(\mu\)-mass is
\(O(B^{-0.49})\).  The medium-prime deletion has
\begin{equation}\label{eq:medium-drop}
 \mathbb E\left[W_{\rm drop}(F_{\rm med})
  +\sum_{\mathfrak q\in F^\circ}p^{-\alpha_s}
                    \Delta_{\mathfrak q}(F_{\rm med})
  +\sum_s\Delta_{\mathfrak l_s}(F_{\rm med})\right]
       \ll B^{-c_0}
\end{equation}
for some fixed \(c_0>0\).  This follows from the rough conditional
bounds and \(|\mathcal Q_{\rm med}|\): the positive logarithmic
costs \(O(L^{b-v}+L^{m_0-b}L_2)\) are smaller than \(\log B\).

Let \(w_i^\circ=(p^{-\alpha_s/2}\log p)\) on the field-\(i\)
labels.  The root probabilities and their two-prime versions give
\[
 \mathbb E\mu(F_i^\circ)=(1+o(1))\sum_{p\in J}p^{-\alpha_s}
       \le(2+o(1))u,
 \quad
 \mathbb E\|w_i^\circ\|^2=(1+o(1))\phi''(\alpha_s),
\]
with negligible variances.  The deterministic protected labels have
negligible weight in these estimates.

Markov's and Chebyshev's inequalities therefore give an event of
probability \(1-o(1)\) on which all of the following hold:
\begin{enumerate}[label=(\roman*)]
\item Both perturbations lie in \(\mathcal B_U\), and
\(|\widehat C-1|\le1/8\).
\item \(\mu(F_i^\circ)\le3u\) and
\(\|w_i^\circ\|^2=(1+o(1))\phi''(\alpha_s)\).
\item \(\mu(D_0)\le\mu_*/2\), and no protected label belongs to
\(D_0\).
\item \(F_A^2,D_A\le E^{74/75}\) and
\(F_\ell^2\le\ell^{-\alpha_s}E^{74/75}\).
\item Each nonnegative term inside the expectation in
\eqref{eq:medium-drop} is at most \(B^{-c_1}\), for a fixed
\(c_1>0\).
\end{enumerate}
For the initial diagonal deletion,
\(\mathbb E\mu(F_{\rm diag})\le\tau^{-2}\mathbb ED_A\ll uE^{1/2}
=o(\mu_*)\).  Protected diagonals are excluded with probability at
most \(d\tau^{-2}E=o(1)\).  Every other expectation-to-threshold
ratio is a fixed positive power of \(E\) times a polynomial in
\(L\), or a fixed negative power of \(B\).

\begin{lemma}[The deletion invariant]\label{lem:peeling}
On the good event, every stage of \eqref{eq:peel-process} satisfies
\(\mu(D_j)\le\mu_*\), and no protected label is deleted.
\end{lemma}

\begin{proof}
Suppose the assertion holds at stage \(j\).  A lost relation hitting
a protected label must hit a medium-prime label or a previously
deleted label over a different prime.  The mean bound and
Cauchy--Schwarz give
\[
 \Delta_{\mathfrak l_s}(D_j)
 \le B^{-c_1}+2\mu_*+\sqrt{\mu_*}\,E^{37/75}=o(\zeta).
\]
Thus no protected label enters \(D_{j+1}\).

For every label newly deleted at any stage up to \(j+1\), its
lost mass exceeded \(\zeta\) before it was deleted.  It is not a
shared-root label, so a previously deleted label over the same
rational prime cannot share its hit.  Subtracting the mean part,
multiplying by \(p^{-\alpha_s}\), and summing over these newly
deleted labels gives, by weighted Cauchy--Schwarz and AM--GM,
\begin{equation}\label{eq:peel-invariant-bound}
 \mu(D_{j+1}\setminus D_0)
 \le4\zeta^{-1}B^{-c_1}
       +4\zeta^{-2}\mu_*F_A^2
 <\frac12\mu_*.
\end{equation}
Indeed, \(\zeta^{-2}E^{74/75}=E^{221/300}\to0\).
Together with \(\mu(D_0)\le\mu_*/2\), this closes the induction.
The same inequality works for the full cumulative newly deleted
set, so no union bound over the number of stages is needed.
\end{proof}

Let \(D_\infty\) be the terminal set, \(F^*=F_{\rm all}\setminus
D_\infty\), and \(\mathcal S^*=\mathcal S(D_\infty)\).
It contains all protected labels and lies in \(F_{\rm clean}\).
Restrict the feature vectors to \(F^*\), and let \(\widetilde H^*\)
be their Gram over \(\mathcal S^*\).
Distinct retained labels with the same underlying prime have
different residues and cannot hit together.  Their centered entries
are \(O(p^{-\alpha_s})\) in blocks of size at most \(2d\), so their
operator norm is \(O(dB^{-0.49})\).  The remaining retained entry
errors are controlled by \(\tau+F_A\).

The total dropped relation weight is \(O(\mu_*)\).  A raw relation
has at most \(5u\) nonzero prime features, while
\(1+\|\boldsymbol\ell(a)\|^2\le L^{O(1)}\).  Applying
Cauchy--Schwarz to each dropped relation and using
\(\Delta_{\mathfrak q}(D_\infty)\le\zeta\) on retained labels gives
\begin{equation}\label{eq:gram-error}
 \|\widetilde H^*-\widehat C\,U^{\rm ex}|_*\|_{\rm op}
 \ll L^{O(1)}\mu_*+O(u\zeta)+\tau+F_A+dB^{-0.49}
 =O(E^{1/10}).
\end{equation}
The bound is independent of the potentially exponential number of
feature rows.

\section{The spectral lower bound after removing the two norm identities}
\label{sec:spectral}

Let \(H\) be the \((d-1)\)-by-\(d\) Helmert matrix: row \(j\)
has \(j\) entries equal to \(1/\sqrt{j(j+1)}\), followed by
\(-j/\sqrt{j(j+1)}\), followed by zeros.  Then
\[
 HH^T=I_{d-1},\qquad H\mathbf1_d=0,\qquad
 H^TH=I_d-d^{-1}\mathbf1_d\mathbf1_d^T.
\]
Replace the target roots by the actual roots to define
\(\boldsymbol\ell_i^{\rm true}(a)), and form the reduced vector
\[
 \mathbf r_{\rm true}(a)=
 (1,(z_{\mathfrak q}(a))_{\mathfrak q\in F^*},
             H\boldsymbol\ell_1^{\rm true}(a),
             H\boldsymbol\ell_2^{\rm true}(a)).
\]
Its Gram over \(\mathcal S^*\), with weights \(w_{\rm eq}\), is
denoted \(H_{\rm true}^*\).

\begin{proposition}[Uniform positive Gram]\label{prop:positive-gram}
On the good event of Section~\ref{sec:peeling},
\[
 \lambda_{\min}(H_{\rm true}^*)
      \ge c_\delta L^{-2h-r}\ge L^{-3}
\]
for all sufficiently large \(N\).
\end{proposition}

\begin{proof}
On the retained field-\(i\) labels put
\(w_i=(p^{-\alpha_s/2}\log p)\).  Deleting mass \(\mu_*\) changes
its squared norm by at most \((\log B)^2\mu_*\), so
\begin{equation}\label{eq:wi-norm}
 \|w_i\|^2\asymp u(\log B)^2\asymp L^{h+b}.
\end{equation}
Every retained relation satisfies the exact identity
\begin{equation}\label{eq:norm-identity}
 \sum_{\mathfrak q\in F_i^*}w_{i,\mathfrak q}z_{\mathfrak q}(a)
      =\log|f_i(a)|
      =\log X_{\rm sc}+\mathbf1_d^T
                               \boldsymbol\ell_i^{\rm true}(a).
\end{equation}
These two identities explain the necessary projection.

Take a unit input vector
\(v=(v_0,v_{F_1},v_{F_2},v_{U_1},v_{U_2})\) in the reduced
feature space.  Put
\begin{align*}
 t_i&=-\frac{w_i^Tv_{F_i}}{\|w_i\|^2},&
 v'_0&=v_0-(t_1+t_2)\log X_{\rm sc},\\
 v'_{F_i}&=v_{F_i}+t_iw_i,&
 v'_{\rm arch,i}&=H^Tv_{U_i}-t_i\mathbf1_d.
\end{align*}
Then \(v'_{F_i}\perp w_i\), and \eqref{eq:norm-identity} shows
that the corresponding linear forms on true features agree exactly.
Moreover,
\[
 \|v'\|^2=(v'_0)^2+
   \sum_i\bigl(\|v_{F_i}\|^2-t_i^2\|w_i\|^2
                       +\|v_{U_i}\|^2+dt_i^2\bigr).
\]
If a substantial fraction of the original norm lies in the removed
prime directions, the last terms give \(cd/\max_i\|w_i\|^2\).
Otherwise either the remaining components have substantial norm,
or cancellation of a substantial \(v_0\) forces
\(d(t_1^2+t_2^2)\gg d/(\log X_{\rm sc})^2\).  Thus
\begin{equation}\label{eq:projection-loss}
 \|v'\|^2\ge c\min\left(1,
        \frac d{\max_i\|w_i\|^2},
        \frac d{(\log X_{\rm sc})^2}\right)
       \gg L^{a-2h}.
\end{equation}

Normalize \(y=v'/\|v'\|\).  Let \(U(a)\) be the covariance model
obtained by replacing the exact same-field second moments with the
first two terms of \eqref{eq:covariance}.  The total Frobenius error
between \(U^{\rm ex}\) and \(U\) is \(O(u^{-1})\), because
\(\mu(F_i^*)=O(u)\).  Crucially, this error is supported only on
the prime-feature block.  It is absorbed there by the order-one
positive diagonal; applying a global Weyl estimate against the much
smaller final gap would not be valid.

The negative covariance direction in field \(i\) differs from
\(w_i/\sqrt{\phi''(\alpha_{\mu_i(a)})}\) in norm by \(O(u^{-1})\).
Since \(y_{F_i}\perp w_i\), its quadratic contribution is
\(O(u^{-2})\|y_{F_i}\|^2\).  Also
\(d_{\mathfrak q}-m_{\mathfrak q}^2=1+O(u^{-1})\), and
\(\sum_{\mathfrak q}m_{\mathfrak q}(a)^2\le12u\).
Completing the square, with a small multiple of \(1/u\), therefore
gives
\begin{equation}\label{eq:model-coercive}
 y^TU^{\rm ex}|_*y
 \ge\frac12\|y_F\|^2+
       \frac c u\sum_a\pi_0(a)
             (y_0+\boldsymbol\ell(a)^Ty_{\rm arch})^2.
\end{equation}

It remains to control the archimedean part.  Let
\(\overline{\boldsymbol\ell}^{(j)}\) be its average on
\(I_j\cap\mathcal A^*\).  The \(2d\)-by-\(2d\) matrix with
rows \(\overline{\boldsymbol\ell}^{(j)}-
\overline{\boldsymbol\ell}^{(0)}\), \(1\le j\le2d\), has
diagonal magnitude \(\kappa L^t+O(1)\) and off-diagonal Frobenius
norm \(O(d\log d)\).  Thus its smallest singular value is
\(\gg L^t\), since \(t>a\).
The reference average has squared norm \(O(d)\), since its
interval is bounded away from all target roots.
Writing \(s_j=y_0+
\overline{\boldsymbol\ell}^{(j)T}y_{\rm arch}\), the singular-value
bound gives
\(\sum_j s_j^2\gg L^{2t}d^{-1}\|y_{\rm arch}\|^2\).
Together with
\(y_0^2\le2s_0^2+O(d)\|y_{\rm arch}\|^2\) and \(t>a\), this yields
\[
 \sum_a\pi_0(a)(y_0+\boldsymbol\ell(a)^Ty_{\rm arch})^2
       \ge\frac c d(y_0^2+\|y_{\rm arch}\|^2).
\]
Equation~\eqref{eq:model-coercive} is consequently at least
\(c/(du)\).

The error \eqref{eq:gram-error} is smaller than every inverse power
of \(L\).  Replacing target roots by actual roots costs at most
\(L^{O(1)}K_1^{-1.8}\) in the squared linear-form comparison.
Since \(\widehat C\ge7/8\), these errors can be absorbed.  Combining
with \eqref{eq:projection-loss} gives
\[
 v^TH_{\rm true}^*v\gg
        \frac{L^{a-2h}}{du}\asymp L^{-2h-r}.
\]
Finally \(2h+r=5/3+21\delta<3\), proving the claimed guard.
\end{proof}

\begin{corollary}[A successful enumerated pair]\label{cor:pair-exists}
For every sufficiently large surviving input and every
\(\ell\in\mathcal P_{\rm gen}^{\rm good}\), there is a pair
\((\mathbf u,\mathbf v)\in\mathcal B_U^2\) with a retained set
\(F^*\subseteq F_{\rm clean}\), containing all protected labels,
whose reduced true Gram is positive definite.
\end{corollary}

\begin{proof}
The good event has probability \(1-o(1)\), hence at least \(1/2\)
for sufficiently large \(N\), and it includes membership in the
finite box.  Positive probability proves existence in that box.
The deterministic search enumerates every pair in it.
\end{proof}

\section{Exact order extraction from the two real-unit lattices}
\label{sec:extraction}

This section records the finite algebraic interface used by the analytic
proof.  These statements are parameter-general.  Their hypotheses do
not contain the original decimal powers of \(L\).

\subsection{The nonmonic valuation offset}

Let \(f(X)=\sum_{j=0}^dc_jX^j\) be primitive and irreducible,
\(K=\mathbb Q(\theta)\), and \(f(\theta)=0\).  For a maximal ideal
\(\mathfrak p\) of \(\mathcal O_K\), use the integer-normalized
valuation \(v_{\mathfrak p}\).  Define
\[
 w_{\mathfrak p}(a,b)=v_{\mathfrak p}(a-b\theta)
                       -\min(0,v_{\mathfrak p}(\theta)).
\]
For \((a,b)=1\), source Lemma~7.3 gives
\begin{equation}\label{eq:valuation-sum}
 w_{\mathfrak p}(a,b)\ge0,\qquad
 \sum_{\mathfrak p\mid p}f(\mathfrak p/p)
                 w_{\mathfrak p}(a,b)=v_p(F_f(a,b)),
\end{equation}
where \(F_f(a,b)=b^df(a/b)\), interpreted homogeneously at \(b=0\).
To verify the offset, the local factorization of \(f\) over
\(\mathbb Q_p\) and the Gauss valuation give
\[
 \sum_{\mathfrak p\mid p}f(\mathfrak p/p)
                 \min(0,v_{\mathfrak p}(\theta))=-v_p(c_d).
\]
Subtract this from the valuation of
\(N_{K/\mathbb Q}(a-b\theta)=F_f(a,b)/c_d\).
Nonnegativity follows directly from the ultrametric inequality and
\((a,b)=1\).

If \(p\nmid c_d\) and \(r_0\) is a simple root of \(f\bmod p\),
Hensel's lemma gives a unique degree-one, unramified prime ideal
\[
 \mathfrak p_{p,r_0}=p\mathcal O_K+
                    (c_d\theta-c_dr_0)\mathcal O_K.
\]
When \(v_p(f(a))=1\), its valuation feature is exactly
\(\mathbf1_{a\equiv r_0\ (p)}\), and all other ideals above
\(p\) have adjusted valuation zero.  This proof uses the local
simple factor and does not require coprimality of a power-order index
with \(p\).  General ideal and valuation facts are as in
\cite{Neukirch1999}.

\subsection{The augmented valuation matrix and unit lattice}

List all \(a_1<\cdots<a_M\) in
\(\mathcal S(F_{\rm mult}\cup F_{\rm shared})\).  Define the
\(0\)-\(1\) valuation matrix
\[
 V_{\mathfrak q,j}=\mathbf1_{a_j\equiv r_0\ (p)},
 \qquad \mathfrak q=(i,p,r_0)\in F_{\rm clean},
 \qquad
 V_{\rm aug}=\begin{pmatrix}\mathbf1_M^T\\V\end{pmatrix}.
\]
Put \(D_{\rm row}=1+|F_{\rm clean}|\),
\(\Lambda_{\rm val}=\ker(V_{\rm aug})\cap\mathbb Z^M\), and
\(D_U=2d-2\).  The top row is needed to cancel the nonmonic offsets
in \eqref{eq:valuation-sum}.

Let \(U_{\rm arch}\) have columns
\((H\boldsymbol\ell_1^{\rm true}(a_j),
H\boldsymbol\ell_2^{\rm true}(a_j))\), and define
\[
 \Lambda_U=
 \{(HL_1(\epsilon_1),HL_2(\epsilon_2)):
                      \epsilon_i\in\mathcal O_{K_i}^{\times}\},
\]
where \(L_i\) is the vector of real-embedding logarithms.  By
Dirichlet's unit theorem this is a full-rank lattice in
\(\mathbb R^{D_U}\).  Its covolume has the elementary lower bound
\begin{equation}\label{eq:unit-covolume}
 \det\Lambda_U>(9d)^{-d}.
\end{equation}
Indeed, if a totally real unit \(\epsilon\ne\pm1\) had all
\(|\sigma(\epsilon)|^2<2\), every factor of
\(|N(\epsilon^2-1)|\) would be smaller than one, contradicting
its positive integral norm.  Thus its log vector has norm greater
than \(1/3\).  Minkowski's convex-body theorem applied to the cube
of side \(2/(3\sqrt{d-1})\) gives the claimed product-lattice
covolume.  The only units with zero log vector are \(\pm1\).

Let \(b_\ell\) be the valuation-row vector equal to one on the
\(d\) protected labels and zero elsewhere, including the top row.
If
\[
 V_{\rm aug}c=-E b_\ell,\qquad c\in\mathbb Z^M,
\]
then \eqref{eq:valuation-sum} and the top row show that
\[
 \gamma_1=\ell^E\prod_j(a_j-\theta_1)^{c_j},\qquad
 \gamma_2=\prod_j(a_j-\theta_2)^{c_j}
\]
are units.  In field one, the protected primes exhaust all ideals
above \(\ell\), because the polynomial splits completely and
simply there.  Their valuations are canceled by \(\ell^E\).
All other valuations vanish.  Their projected log vector is
\(U_{\rm arch}c\).

\begin{lemma}[Exact two-field certificate]\label{lem:unit-certificate}
Suppose there are \(c^{(0)}\in\mathbb Z^M\), \(E_0\ge1\), and
\(c^{(1)},\ldots,c^{(D_U)}\in\Lambda_{\rm val}\) such that
\(V_{\rm aug}c^{(0)}=-E_0b_\ell\), and the matrix
\[
 M_U=(U_{\rm arch}c^{(1)},\ldots,U_{\rm arch}c^{(D_U)})
\]
is nonsingular.  Then
\(x=-M_U^{-1}U_{\rm arch}c^{(0)}\) is rational.  If \(Q\) is
the least common multiple of its reduced denominators, then
\begin{equation}\label{eq:exact-order}
 \ell^{2E_0Q}\equiv1\pmod N.
\end{equation}
\end{lemma}

\begin{proof}
All columns of \(M_U\), and \(U_{\rm arch}c^{(0)}\), lie in
\(\Lambda_U\).  The sublattice \(M_U\mathbb Z^{D_U}\) has
positive integer index
\(E_1=|\det M_U|/\det\Lambda_U\).  Thus \(E_1x\in\mathbb Z^{D_U}\),
and every denominator of \(x\), hence \(Q\), divides \(E_1\).
Set
\[
 c=Qc^{(0)}+\sum_{\nu=1}^{D_U}(Qx_\nu)c^{(\nu)}.
\]
It is integral, has valuation vector \(-E_0Qb_\ell\), and has
zero projected unit logarithms.  Both resulting units are \(\pm1\).

Evaluation \(\theta_i\mapsto m\pmod N\) defines a homomorphism
on \(\mathbb Z[c_d^{-1}][\theta_i]\), since the defining polynomial
is monic over \(\mathbb Z[c_d^{-1}]\) and vanishes at \(m\bmod N\).
Localize further at the finitely many \(a_j-\theta_i\).  Their
images \(a_j-m\) are invertible by the initial gcd checks, so the
homomorphism extends.  The two unit products above belong to this
localized ring.  Their images differ by \(\ell^{E_0Q}\) and are
both signs.  Dividing and squaring proves \eqref{eq:exact-order}.
No homomorphism from the entire maximal order is assumed.
\end{proof}

Positive definiteness in Corollary~\ref{cor:pair-exists} gives exactly
the required spanning hypotheses.  The retained feature columns span
their whole real space; their valuation parts therefore contain
\(b_\ell\), and the image of the real valuation kernel under
\(U_{\rm arch}\) is all of \(\mathbb R^{D_U}\).
The valuation matrix is integral, so its real kernel is spanned by
its rational kernel and hence by \(\Lambda_{\rm val}\).
Extending combinations by zero to the omitted clean relations
preserves these conclusions for the full matrices.

\subsection{The exact finite extractor interface}

The following is source Lemma~7.5, with its input promises stated
explicitly.  It is the deterministic procedure used in
Algorithm~\ref{alg:annihilator}.

\begin{lemma}[\(\operatorname{TryPair}\), source Lemma 7.5]
\label{lem:trypair}
Let \(N\ge9\), \(m\ge2\), \(d\ge9\), \(c_d\ge1\),
\(A\ge2\), \(B\ge4\), \(K_1\ge16\), and \(L_0\ge10\) be
integers, with
\[
 N\le2^{dL_0},\quad(c_d,N)=1,\quad
 m,A,B,K_1\le2^{L_0}.
\]
Let \(J\) be a duplicate-free prime list at most \(B\), coprime
to \(c_dN\), and let \(\mathcal A^*\subseteq[1,A]\cap\mathbb Z\)
be a duplicate-free list with \((a-m,N)=1\).
Let \(f_1,f_2\) be primitive, irreducible degree-\(d\) integer
polynomials with common leading coefficient \(c_d\),
\(f_i(m)\equiv0\pmod N\), and height at most \(2^{L_0}\).
Put
\[
 P_{\rm bit}=64d^3D_{\rm row}L_0,
 \qquad e_{\rm rt}=\lceil(19/20)\lfloor\log_2K_1\rfloor\rceil,
 \qquad 2^{-e_{\rm rt}}<1/(8d).
\]
There is a deterministic procedure, polynomial in
\(A,B,d,D_{\rm row},P_{\rm bit}\), which returns either failure
or a verified even \(E_\ell\ge2\) with
\(\log_2E_\ell\le P_{\rm bit}\) and \(\ell^{E_\ell}=1\pmod N\).
It succeeds whenever:
\begin{enumerate}[label=(\roman*)]
\item The protected split-root hypothesis holds and all protected
labels are clean.
\item Each \(f_i\) has opposite signs at
\(A(\rho_{i,s}\pm2^{-e_{\rm rt}})\), for all \(s\).
\item Every clean relation has at most \(L_0\) valuation features,
including the top row, and distance at least \(2^{-L_0}\) from
every normalized real root.
\item \(b_\ell\in\operatorname{im}_{\mathbb Q}(V_{\rm aug})\) and
\(U_{\rm arch}(\Lambda_{\rm val})\) spans \(\mathbb R^{D_U}\).
\end{enumerate}
\end{lemma}

For reconstruction, use the five steps specified in the source's
proof of Lemma~7.5: enumerate and factor the bounded values by trial
division over \(J\); build the integer valuation matrix and its
exact pivot minors; compute an integral valuation solution and
integral kernel vectors; isolate the real roots and compute certified
logarithmic Gram data; and use continued fractions to reconstruct
the rational coordinates in Lemma~\ref{lem:unit-certificate}.
The last step verifies the resulting modular power directly.

We explain why this is a finite precision statement.  Integer
pivot minors are bounded by Hadamard's inequality using the
\(L_0\) nonzero entries per column.  Thus the kernel coefficients
have bit lengths polynomial in \(D_{\rm row}L_0\).
The archimedean columns have bounded logarithms because of the
root-distance guard.  Equation~\eqref{eq:unit-covolume} bounds
the index denominators and provides a separation bound for
nonsingular logarithmic minors.  Exact bordered-adjugate updates
select the required independent columns, and interval arithmetic
at a polynomial multiple of \(P_{\rm bit}\) bits is enough for
the Gram tests and rational reconstruction.  The proof of source
Lemma~7.5 supplies the explicit constants, including the displayed
\(64d^3D_{\rm row}L_0\) cap.  Its hypotheses have been retained
rather than replaced by a heuristic numerical rank test.

By Lemma~\ref{lem:construction}, all input promises and conditions
(i)--(iii) hold for sufficiently large inputs and every enumerated
pair.  The sign argument used there also holds at the possibly
factor-two smaller radius \(2^{-e_{\rm rt}}\), since
\(K_1^{0.05}\) dominates \((8d)^d\).
Condition (iv) follows from a successful Gram pair.  The positive
Gram and the proof-only deleted set need not be computed by
\(\operatorname{TryPair}\); they guarantee that its exact algebraic
search on the full clean relation set succeeds.

\section{From many certified primes to a common exponent}
\label{sec:saturation}

The saturation step is elementary but important.  Certificates for
many individual primes do not automatically annihilate every unit
modulo every unknown prime factor.

\begin{lemma}[Subgroup-index saturation]\label{lem:saturation}
Suppose \(P^-(N)>K\ge B\),
\[
 K\ge B\exp\left(\frac L{\log B}\log(6L)\right),
\]
and \(\mathcal P\subset(B/2,B]\) is a prime set with
\(|\mathcal P|\ge B/(6\log B)\).  Suppose
\(\ell^{M_*}=1\pmod N\) for every \(\ell\in\mathcal P\).
Let
\[
 L_{\rm idx}=\prod_{p\le K}p^{\max\{j:p^j\le K\}}
              =\operatorname{lcm}(1,2,\ldots,K),
 \qquad M=M_*L_{\rm idx}.
\]
Then \(p-1\mid M\) for every prime \(p\mid N\), and
\(\log L_{\rm idx}\le2K\log4\) for \(K\ge20\).
\end{lemma}

\begin{proof}
Fix \(p\mid N\), and let \(G\) be the subgroup of
\(\mathbb F_p^\times\) generated by the reductions of
\(\mathcal P\).  Its order divides \(M_*\).  Put
\(j=\lfloor\log(p-1)/\log B\rfloor\ge1\).
Products of \(j\)-element multisets from \(\mathcal P\) are
distinct positive integers at most \(B^j\le p-1\), hence distinct
elements of \(G\).  Unique factorization and the elementary
binomial bound give
\[
 |G|\ge {|\mathcal P|+j-1\choose j}
       \ge\left(\frac{|\mathcal P|}{j}\right)^j
       \ge\left(\frac B{6L}\right)^j.
\]
Since \(p-1<B^{j+1}\), the index \(I=[\mathbb F_p^\times:G]\)
satisfies
\[
 I<B(6L)^j\le B\exp\left(\frac L{\log B}\log(6L)\right)\le K.
\]
Every prime power dividing \(I\) is therefore at most \(K\), so
\(I\mid L_{\rm idx}\).  Thus \(p-1=|G|I\mid M\).
The stated logarithmic bound follows from
\(\log L_{\rm idx}=\sum_{j\ge1}\vartheta(K^{1/j})\) and the
elementary Chebyshev bound \(\vartheta(x)\le x\log4\).
\end{proof}

The rational guard in Algorithm~\ref{alg:annihilator} implies
\(|\mathcal P_{\rm ok}|\ge B/(6\log B)\), because
\(e(b)\kappa_-<\log B\).  Every recorded prime has a verified
order certificate, so every element of the generated subgroup is
annihilated by their least common multiple.  Thus, whenever the
algorithm returns \(M\), its correctness follows from
Lemma~\ref{lem:saturation} on every input, independently of the
asymptotic existence proof.

\section{The supplied-exponent splitter}
\label{sec:splitter}

We only need the original, slower splitter.  The companion
order-splitting manuscript proves its finite algebraic and analytic
interfaces, and also studies a faster family.  The following weaker
statement suffices here.

\begin{theorem}[Supplied-exponent splitting]\label{thm:splitter}
There is an effective deterministic algorithm which, given a
composite non-perfect-power integer \(N\) and an even integer
\(M\ge2\) satisfying \(p-1\mid M\) for all \(p\mid N\),
returns a proper factor in
\[
 \exp(O(\sqrt L\log\log(N+16)))\,(log M)L^2
\]
bit operations.  The factorization of \(M\) is not required.
\end{theorem}

For orientation, write \(M=2^tu\), \(u\) odd.  Testing the successive
squares of \(z^u\) and gcds with \(N\) either splits \(N\) or
forces the same exact power-of-two order at every prime-power
component.  One unknown prime or product of two unknown primes
then defines an active quadratic character whose value \(-1\)
would force a split.  Failure on tested primes produces compatible
modular square roots in quadratic and biquadratic fields.  Enumerating
unit and small-principal-ideal generators forces a positive ideal
character family.  Its Walsh positivity contradicts a uniform
quadratic-field Burgess bound at an effectively reached dyadic scale.

The exact proof dependencies are source Lemma~3.1, the corrected
enumeration contract of Lemma~4.2, Lemma~5.1, the corrected
finite-field character-sum input to Lemma~6.3(1), and the
Lemma~6.4/Theorem~1.2 search.  The companion proof corrects two
literal statements: the squarefactor in the root construction need
only divide a power of the primorial, and the general finite-field
Weil statement counts distinct roots in an algebraic closure.
The actual moment polynomials split into linear factors, so that
correction preserves the application.  The enumeration replacement
uses exact quadratic units and a dimension-four weighted lattice
search, with polynomial dependence on the numerical discriminant.
This is a proved modular input to the present factoring argument.

\section{Correctness, bit complexity, and effective finite-input coverage}
\label{sec:runtime}

We now assemble the proof of Theorem~\ref{thm:main}.
Every direct factor output is checked.  Every returned \(E_\ell\)
is checked by modular exponentiation.  The exact density guard and
Lemma~\ref{lem:saturation} certify every returned common exponent
\(M\).  Applying Theorem~\ref{thm:splitter} therefore returns a
factor whenever the annihilator returns an exponent.  Complete
trial division guarantees termination on all other inputs.

For sufficiently large surviving inputs, all construction guards
hold by \eqref{eq:gates1}--\eqref{eq:gates3} and
Lemma~\ref{lem:construction}.  By
Corollary~\ref{cor:pair-exists} and Lemma~\ref{lem:trypair}, every
\(\ell\in\mathcal P_{\rm gen}^{\rm good}\) succeeds somewhere in
the enumerated box.  Lemma~\ref{lem:exception-count} gives more
than enough such primes to pass the exact density guard.

Since \(|F_{\rm clean}|\le2dB\), each certificate has
\[
 \log_2E_\ell\le64d^3(1+2dB)L_0.
\]
There are at most \(B\) certificates, and taking an lcm cannot
exceed the bit length of their product.  Hence
\begin{equation}\label{eq:exponent-size}
 \log M\le L^{O(1)}B^2+2K_{\rm idx}\log4
       =\exp(O(L^{1-b}L_2)).
\end{equation}
Because \(g=1-b+\delta\), this is smaller than the bit-length
cap \(2^{e(g)}\) eventually.  No large input invokes fallback.

The total number of extractor calls is at most
\begin{equation}\label{eq:call-count}
 |\mathcal P_{\rm gen}|\,|\mathcal B_U|^2
 \le16BL_0^{80}U^4A^2
       =\exp((4\kappa+o(1))L^h).
\end{equation}
Each call has bit complexity
\(\operatorname{poly}(A,B,d,D_{\rm row},P_{\rm bit})
=\exp(O(L^c))\), which is a lower power because \(c<h\).
All base constructions are polynomial in \(B\).  Trial division,
the gcd loop, and saturation cost
\(\operatorname{poly}(K_{\rm idx},A,L)
=\exp(O(L^{1-b}L_2+L^c))\), again lower powers.
Finally, \eqref{eq:exponent-size} and
Theorem~\ref{thm:splitter} give splitting cost
\[
 \exp(O(\sqrt L L_2+L^{1-b}L_2))=\exp(o(L^h)).
\]
This proves the stated \(\exp(O_\delta(L^h))\) bound.
In fact, the leading enumeration coefficient in
\eqref{eq:call-count} is \(4\log2<3\), so an effective
\(C_\delta\exp(3L^h)\) upper bound also follows after absorbing
the finite initial range.

\subsection{Why the threshold is effective}

All computational parameters are rationally rounded.  The algorithm
uses no zero-finding subroutine and no test for an exceptional
Dirichlet character.  The zero sets classify the analysis of sample
points; they do not define the search.  The exceptional real
character is handled by its permitted nonpositive prime-sum main
term, without an ineffective estimate for the distance of its zero
from one.

Every asymptotic comparison in the proof has an explicit strict
power gap, or a fixed positive power of \(E\).  Effective constants
in the stated classical inputs therefore give an effective
threshold \(N_0(\delta)\).  Above it, all finite guards hold and
enough searches succeed.  Below it, trial division has a finite
maximum cost, absorbed by an effective constant depending on
\(\delta\).  The executable algorithm does not need
\(N_0(\delta)\) as an input or precomputed datum.

For a requested fixed \(\varepsilon>0\), choose any fixed rational
\(0<\delta\le1/100\) with \(8\delta<\varepsilon\).
Then \(h<2/3+\varepsilon\), completing the proof.

\section{What currently prevents a further power reduction}
\label{sec:barrier}

The proof identifies a concrete barrier rather than treating
smoothness equidistribution as a black box.  The main progression
argument needs two prime prefixes larger than its modulus.  This
requires \(q^2B^2\ll Y\), and the actual moduli include a factor
\(|a-m|\) of size approximately \(m\).

Suppose the two prefix thresholds are shortened to
\(Z_1,Z_2\le q\).  For a fixed terminal prime, orthogonality only
gives
\[
 \sum_{\chi\bmod q}|A_i|^2
       \ll\varphi(q)Z_i^{1-2\alpha}.
\]
Cauchy--Schwarz leaves
\begin{equation}\label{eq:short-prefix-loss}
 \frac q{\varphi(q)}\sum_\chi|A_1A_2|
       \ll q(Z_1Z_2)^{1/2-\alpha}.
\end{equation}
At \(\alpha=1-o(1)\), avoiding a loss exponential in \(\log q\)
requires \(Z_1Z_2\ge q^{2-o(1)}\).  If
\(q\ge Y^{1/2+\epsilon_0}\) for a fixed \(\epsilon_0>0\), even
the maximal available product \(Z_1Z_2=O(Y)\) leaves a loss of
\(\exp((\epsilon_0+o(1))\log Y)\).  The Euler-product saving is
only \(\exp(-O(u_Y))\), where
\(u_Y=\log Y/\log B=o(\log Y)\).  Smoothing the window cannot
alter this arithmetic size comparison.

This is a limitation of the orthogonality/Cauchy--Schwarz proof,
not a lower bound on factoring or on equidistribution.  An average
over the special moduli arising from the construction, an estimate
that exploits their actual residue classes, or a stronger
multilinear character bound might behave differently.

The present height mechanism has a second transparent restriction.
If \(d=L^a\), \(\log A\asymp L^c\), and the perturbation width
dominates both \(m\) and \(A^d\), then
\[
 h\ge1-a,\qquad h\ge a+c,
\]
while root separation and the smoothness/sieve layers require
\(c>b>k>t>a\).  Thus \(h>2a\), and minimizing
\(\max(1-a,2a)\) gives the limiting value \(2/3\) at
\(a=1/3\).  This optimization describes the mechanism proved here.

\subsection{Why translating all roots towards the evaluation point fails}

A tempting way to shrink \(|a-m|\) is to place all polynomial roots
near \(m\).  The exact product formula shows the cost.  Suppose
\(f(m)=kN\ne0\), every real root satisfies
\(|m-\theta_j|\le CA\), and a sampled \(a\) is at distance at
least \(wA\) from its nearest root and at least
\(A/(C'd)\) from all the others.  Then
\[
 \frac{|f(a)|}{|f(m)|}
   =\prod_j\frac{|a-\theta_j|}{|m-\theta_j|}
   \ge wC^{-d}(C'd)^{-(d-1)}.
\]
Consequently
\(\log|f(a)|\ge\log N-\log(1/w)-O(d\log d)\).
In the intended regime these norms are \(N^{1-o(1)}\), losing
the subexponential value size.  This rules out a simple translation
of the present separated-root construction; it leaves other
polynomial shapes and different relation constructors open.

\subsection{Two exact refinements that may help an endpoint analysis}

For the equally spaced roots used here,
\[
 |P_i'(\rho_{i,s})|
   =\frac{(s-1)!(d-s)!}{(2d)^{d-1}}
   \ge(8e)^{-(d-1)}.
\]
Thus, for \(0<\epsilon\le1/(8d)\),
\[
 |P_i(\rho_{i,s}\pm\epsilon)|>
                     \epsilon\,32^{-(d-1)}.
\]
The coarse root-stability argument can therefore replace a
\(d\log d\) loss by \(O(d)\).  Likewise, grouping the
off-diagonal mean-log entries by root spacing and using
\(\int_0^1\log^2x\,dx=2\) improves their Frobenius bound from
\(O(d\log d)\) to \(O(d)\).  These identities are useful for a
future logarithmic optimization.  They do not by themselves prove
an endpoint \(L_N[2/3,c]\) result: the zero-free height, conductor
exclusions, Gaussian errors, and deletion thresholds would all
have to be rebalanced at logarithmic precision.

\appendix
\section{Source correspondence, corrections, and reconstruction}
\label{app:correspondence}

The proof deliberately separates parameter-general finite algebra
from the scale-sensitive analytic assertions.  The following table
specifies the imported interfaces and the new work needed to use them.

\begin{longtable}{p{.18\textwidth}p{.35\textwidth}p{.37\textwidth}}
\toprule
Source location & Role & Treatment in this manuscript \\
\midrule
\endhead
Lemma 7.3 & Nonmonic ideal valuations & Offset and simple-root
argument reproduced in Section~\ref{sec:extraction}. \\
Lemma 7.4 & Two-field unit cancellation and spanning & Covolume,
index denominators, localization, and Gram implication reproduced. \\
Lemma 7.5 & Finite exact extractor & Full input/output contract
retained in Lemma~\ref{lem:trypair}; its five-step exact implementation
and explicit bit cap are imported. \\
Lemma 8.1 & Base-polynomial CRT construction & Construction interface,
finite recipe, and height transfer given in
Section~\ref{sec:construction}. \\
Lemma 8.2 & Sieve, roots, and means & New power schedule and proof
given in Sections~\ref{sec:parameters}--\ref{sec:sieve}. \\
Lemmas 8.3--8.4 & Gevrey cutoff & Parameter-general kernel and
precise Mellin consequences stated in Section~\ref{sec:progressions}. \\
Lemma 9.1 & Rough progression upper bound & Greedy-prefix proof
reproduced in Lemma~\ref{lem:rankin}. \\
Lemma 9.2 & Prime sums and saddle moments & Horizontal-contour term
retained; Edgeworth and finite-difference proofs supplied. \\
Proposition 9.3 & Clean progressions & Zero-density counts and exact
two-prefix orthogonality argument supplied. \\
Lemmas 9.4--9.5 & Conditional Gaussian laws & Congruence lattices,
coprimality, gcd division, and local exclusions supplied. \\
Proposition 10.1 & Centered Gram and peeling & Centered identity,
conductor costs, adaptive invariant, and blockwise spectral proof
supplied. \\
Lemma 11.1 & Saturation & Elementary subgroup-count proof reproduced. \\
Theorem 1.2(a) & Supplied-exponent splitter & Original slower bound
suffices; complete package treated in the companion manuscript. \\
\bottomrule
\end{longtable}

The corrections used here are narrow and explicit.  The Perron
estimate retains \(\mathcal L/(\nu H)\) on the horizontal segments.
The second-difference proof evaluates four bounded Edgeworth errors
and differentiates only its explicit main term.  The covariance
error is absorbed in the prime-feature block before deriving the
small archimedean gap.  The order-splitting companion corrects the
primorial squarefactor statement and the finite-field root-count
statement.  None of these repairs is concealed inside a numerical
experiment.

To reconstruct the algorithm, use the rounded parameters and
Algorithm~\ref{alg:annihilator}, the frozen source's finite
\(\operatorname{TryPair}\) implementation, and the companion's
supplied-exponent splitter.  To reconstruct the proof, the exact
power-gap checks accompany the repository, while every analytic
comparison needed to justify those gaps is written above.  The
finite checks establish the algebra of the schedule; the mathematical
lemmas establish the uniform probability, spanning, and running-time
claims.
